\subsection{Character of a Representation}

Let E be a left A-module of finite dimension over K. The character of E or trace of E, denoted by \(Tr_{E}\) , is the linear form \(a \mapsto \mathrm{Tr}(a_{\mathrm{E}})\) . Let \((e_{1}, \ldots, e_{n})\) be a basis of E and \((e_{1}^{*}, \ldots, e_{n}^{*})\) its dual basis. By definition, we have the relation

\[
\mathrm{Tr} _ {\mathrm{E}} = \sum_ {i = 1} ^ {n} c _ {\mathrm{E}} (e _ {i}, e _ {i} ^ {*}).\tag{7}
\]

The linear form \(Tr_{E}\) belongs to \(\Theta_{\mathrm{E}}(\mathrm{A})\) . We have \(Tr_{E} = Tr_{E'}\) if E and \(E'\) are two isomorphic A-modules. So \(Tr_{E}\) depends only on the isomorphism class of E.

Let \(\pi\) be a linear representation of A in a vector space V of finite dimension over K. The trace of \(\pi\) , or sometimes character of \(\pi\) , is defined as the character of the A-module \(V_{\pi}\) . We denote it by \(\mathrm{Tr}_{\pi}\) . If \((e_1, \ldots, e_n)\) is a basis of V over K and if \((\pi_{ij}(a))\) is the matrix of \(\pi(a)\) with respect to this basis, then we have

\[
\operatorname{Tr} _ {\pi} (a) = \sum_ {i = 1} ^ {n} \pi_ {i i} (a).\tag{8}
\]

The linear form \(\mathrm{Tr}_{\pi}\) belongs to \(\Theta_{\pi}(\mathrm{A})\) .

PROPOSITION 5. --- Let S be a simple A-module of finite dimension over K, let D be the opposite field of the commutant of S, and let Z be the center of D. The following properties are equivalent:

\begin{enumerate}
\def\labelenumi{(\roman{enumi})}
\item
  The character \(Tr_{S}\) of S is not zero.
\item
  There exists an element \(d \in \mathrm{D}\) such that \(\operatorname{Tr}_{\mathrm{D} / \mathrm{K}}(d) \neq 0\) .
\item
  The extension \(\mathbf{Z}\) of \(\mathbf{K}\) is separable, and the characteristic \(p\) of \(\mathbf{K}\) does not divide the degree \([\mathrm{D}:\mathrm{Z}]\) .
\end{enumerate}

The right D-vector space S is finite-dimensional. Let \((e_{1},\ldots,e_{n})\) be a basis of S over D, and let u be an element of \(\operatorname{End}_{\mathrm{D}}(\mathrm{S})\) . Let \((d_{ij})\) be the matrix of u with respect to the basis \((e_{1},\ldots,e_{n})\) . We have \(u(e_{j})=\sum_{i=1}^{n}e_{i}d_{ij}\) for \(j\in[1,n]\) . Denote by \(u_{K}\) the mapping u viewed as an endomorphism of the K-vector space S and by \((u_{ij})\) the matrix of \(u_{K}\) with respect to the decomposition \((e_{i}\mathrm{D})\) of the K-vector space S as a direct sum (II, §10, No.~5, p.~346). We have \(\operatorname{Tr}(u_{\mathrm{K}})=\sum_{i}\operatorname{Tr}(u_{ii})\) . Moreover, \(u_{ii}\) is the endomorphism of the K-vector space \(e_{i}D\) defined by \(u_{ii}(e_{i}d)=e_{i}d_{ii}d\) for \(d\in D\) , so its trace is equal to \(\operatorname{Tr}_{\mathrm{D}/\mathrm{K}}(d_{ii})\) . We have thus proved the equality

\[
\mathrm{Tr} (u _ {\mathrm{K}}) = \mathrm{Tr} _ {\mathrm{D/K}} \bigl (\sum_ {i} d _ {i i} \bigr).\tag{9}
\]

By Burnside's theorem (VIII, p.~83, Corollary 1 of Proposition 4), the mapping \(a \mapsto a_{\mathrm{S}}\) from A to \(\mathrm{End}_{\mathrm{D}}(\mathrm{S})\) is surjective. The equivalence of properties (i) and (ii) therefore follows from formula (9).

If the extension Z of K is separable, then there exists an element \(d \in Z\) such that \(\operatorname{Tr}_{\mathrm{Z/K}}(d) \neq 0\) (V, §8, No.~2, p.~49, Corollary of Proposition 1). Moreover, we have \(\operatorname{Tr}_{\mathrm{D/K}}(d) = [\mathrm{D} : \mathrm{Z}] \operatorname{Tr}_{\mathrm{Z/K}}(d)\) (III, §9, No.~4, p.~548, Corollary); consequently, if p does not divide {[}D : Z{]}, then we have \(\operatorname{Tr}_{\mathrm{D/K}}(d) \neq 0\) . This proves that (iii) implies (ii).

If the extension \(\mathrm{Z}\) of \(\mathrm{K}\) is not separable, then we have \(\operatorname{Tr}_{\mathrm{Z / K}}(x) = 0\) for every \(x \in \mathrm{Z}\) , and therefore \(\operatorname{Tr}_{\mathrm{D / K}}(d) = \operatorname{Tr}_{\mathrm{Z / K}}(\operatorname{Tr}_{\mathrm{D / Z}}(d)) = 0\) for every \(d \in \mathrm{D}\) .

Now suppose that p divides the degree \([D:Z]\) . It then divides the reduced degree of D over Z. For every \(d \in D\) , we have \(\mathrm{Tr}_{\mathrm{D}/\mathrm{Z}}(d) = 0\) (VIII, p.~344, Remark), and therefore \(\mathrm{Tr}_{\mathrm{D}/\mathrm{K}}(d) = \mathrm{Tr}_{\mathrm{Z}/\mathrm{K}}(\mathrm{Tr}_{\mathrm{D}/\mathrm{Z}}(d)) = 0\) . This completes the proof of the implication (ii) \(\Rightarrow\) (iii).

COROLLARY. --- If the field \(\mathbf{K}\) is perfect, then properties (i) through (iii) hold.

Since the field is perfect, the extension Z of K is separable (V, §15, No.~5, p.~125, Theorem 3). Property (iii) then follows from Corollary 3 of VIII, p.~323.

PROPOSITION 6. --- Let \(\mathcal{S}_0\) be the set of classes of simple A-modules of finite dimension over K with nonzero trace. The family of linear forms \((\mathrm{Tr}_{\mathrm{S}})_{\mathrm{S}\in \mathcal{S}_0}\) is free over K.

Let \(\mathrm{F}\) be a finite subset of \(\mathcal{S}_0\) , and let \(\mathrm{S}\) be an element of \(\mathrm{F}\) . By assumption, there exists an element \(a \in \mathrm{A}\) such that \(\operatorname{Tr}_{\mathrm{S}}(a) \neq 0\) . By Corollary 1 of Proposition 4 (VIII, p.~83), there exists an element \(b \in \mathrm{A}\) such that \(b_{\mathrm{S}} = a_{\mathrm{S}}\) and \(b_{\mathrm{T}} = 0\) for every \(\mathrm{T} \in \mathrm{F} - \{\mathrm{S}\}\) . We have \(\operatorname{Tr}_{\mathrm{S}}(b) \neq 0\) and \(\operatorname{Tr}_{\mathrm{T}}(b) = 0\) for \(\mathrm{T} \in \mathrm{F} - \{\mathrm{S}\}\) . The family \((\operatorname{Tr}_{\mathrm{S}})_{\mathrm{S} \in \mathrm{F}}\) is therefore free. Proposition 6 follows.

Remark. --- Let \(S_{K}\) be the set of classes of simple A-modules of finite dimension over K. Proposition 6 also follows from the fact that the sum of the \(\Theta_{\mathrm{S}}(\mathrm{A})\) , for S running through \(S_{K}\) , is direct.

Let

\[
0 \longrightarrow \mathrm{E} ^ {\prime} \longrightarrow \mathrm{E} \longrightarrow \mathrm{E} ^ {\prime \prime} \longrightarrow 0
\]

be an exact sequence of A-modules, all finite-dimensional over K. By Proposition 1 of III, §9, No.~2, p.~543, we have \(\mathrm{Tr}_{\mathrm{E}} = \mathrm{Tr}_{\mathrm{E}'} + \mathrm{Tr}_{\mathrm{E}''}\) . By the definition of the Grothendieck group \(\mathrm{R}_{\mathrm{K}}(\mathrm{A})\) (VIII, p.~194) and its universal property (VIII, p.~186, Proposition 4), there exists a homomorphism of additive groups \(\theta\) from \(\mathrm{R}_{\mathrm{K}}(\mathrm{A})\) to \(\mathrm{A}^*\) characterized by the relation \(\theta([\mathrm{E}]) = \mathrm{Tr}_{\mathrm{E}}\) for every A-module E of finite dimension over K. In particular, the trace of a semisimplification of E is equal to that of E. We deduce from \(\theta\) a K-linear mapping \(\theta_{\mathrm{K}}: \mathrm{K} \otimes_{\mathbf{Z}} \mathrm{R}_{\mathrm{K}}(\mathrm{A}) \to \mathrm{A}^*\) .

COROLLARY. --- a) Suppose that the field K has characteristic 0. The homomorphisms \(\theta\) and \(\theta_{K}\) are injective. Two semisimple A-modules of finite dimension over K are isomorphic if and only if their characters are equal.

\begin{enumerate}
\def\labelenumi{\alph{enumi})}
\setcounter{enumi}{1}
\tightlist
\item
  Suppose that the field K is perfect of nonzero characteristic p.~Then the homomorphism \(\theta_{K}\) is injective, and the kernel of \(\theta\) is \(p\mathrm{R}_{\mathrm{K}}(\mathrm{A})\) . Let E be a semisimple A-module of finite dimension over K. The linear form \(Tr_{E}\) is zero if and only if there exists an A-module F such that E is isomorphic to \(F^{p}\) .
\end{enumerate}

Let \(S_{K}\) be the set of classes of simple A-modules of finite dimension over K. The elements {[}S{]}, where S runs through \(S_{K}\) , form a basis of the Z-module \(\mathrm{R}_{\mathrm{K}}(\mathrm{A})\) , so the elements \(1 \otimes [S]\) form a basis of the K-vector space \(\mathrm{K} \otimes_{\mathbf{Z}} \mathrm{R}_{\mathrm{K}}(\mathrm{A})\) (VIII, p.~195).

Suppose that the field K is perfect. By Proposition 6 and the corollary of VIII, p.~383, the elements \(\theta([S]) = \theta_{\mathrm{K}}(1 \otimes [S]) = \mathrm{Tr}_{\mathrm{S}}\) of \(A^{*}\) are linearly independent over K. It follows that the homomorphism \(\theta_{K}\) is injective and that the kernel of the homomorphism \(\theta\) consists of the elements \(\sum_{S \in S_{K}} n_{S}[S]\) of \(\mathrm{R}_{\mathrm{K}}(\mathrm{A})\) such that \(n_{S} \cdot 1_{K} = 0\) for every \(S \in S_{K}\) . The kernel of \(\theta\) is therefore equal to \(p R_{\mathrm{K}}(\mathrm{A})\) , where p is the characteristic of K. In particular, if K has characteristic 0, then the homomorphism \(\theta\) is injective.

The last assertion of a) follows from the corollary of VIII, p.~190.

Suppose that the assumptions of b) are satisfied, and let E be a semisimple A-module of finite dimension over K. If there exists an A-module F such that E is the direct sum of p submodules isomorphic to F, then the module F is semisimple and finite-dimensional over K and we have \(Tr_{E} = p Tr_{F} = 0\) . Conversely, suppose \(Tr_{E} = 0\) . We have \([E] \in p R_{K}(A)\) , so the multiplicity of every simple module \(S \in S_{K}\) in E is a multiple of p (VIII, p.~190, Proposition 7); we can write it as \(pn_{S}\) with \(n_{S} \in N\) . The family ( \(n_{S}\) ) has finite support. Set \(F = \oplus_{S \in S_{K}} S^{n_{S}}\) . We then have {[}E{]} = {[}F\^{}\{p\}{]}, so that E and F\^{}\{p\} are isomorphic (VIII, p.~190, Corollary).

Let \(\mathrm{E}\) be an A-module of finite dimension over \(\mathrm{K}\) . Let \(a \in \mathrm{A}\) . Denote by \(\chi_{\mathrm{E}}(a; \mathrm{T})\) the determinant of the endomorphism \(1_{\mathrm{E}} + T a_{\mathrm{E}}\) of the \(\mathrm{K}[\mathrm{T}]\) -module \(\mathrm{E}[\mathrm{T}] = \mathrm{K}[\mathrm{T}] \otimes_{\mathrm{K}} \mathrm{E}\) (III, §8, No.~11, p.~541). We have the relation

\[
\chi_ {\mathrm{E}} (a; \mathrm{T}) \equiv 1 + \mathrm{Tr} _ {\mathrm{E}} (a) \mathrm{T} \pmod {\mathrm{T} ^ {2} \mathrm{K} [ \mathrm{T} ]}\tag{10}
\]

(loc. cit., formula (49)). Since this polynomial has constant term equal to 1, it is an invertible formal power series (IV, §4, No.~4, p.~30). Moreover, if

\[
0 \longrightarrow \mathrm{E} ^ {\prime} \longrightarrow \mathrm{E} \longrightarrow \mathrm{E} ^ {\prime \prime} \longrightarrow 0
\]

is an exact sequence of A-modules, all finite-dimensional over K, then we have \(\chi_{\mathrm{E}}(a;\mathrm{T})=\chi_{\mathrm{E}^{\prime}}(a;\mathrm{T})\chi_{\mathrm{E}^{\prime\prime}}(a,\mathrm{T})\) (III, §8, No.~7, p.~534, formula (31)). By the definition of the Grothendieck group \(\mathrm{R}_{\mathrm{K}}(\mathrm{A})\) (VIII, p.~194) and its universal property (VIII, p.~186, Proposition 4), there exists a unique homomorphism \(\chi_{a}\) from the group \(\mathrm{R}_{\mathrm{K}}(\mathrm{A})\) to the multiplicative group \(1+\mathrm{TK}[[\mathrm{T}]]\) such that \(\chi_{a}([\mathrm{E}])=\chi_{\mathrm{E}}(a;\mathrm{T})\) for every A-module E of finite dimension over K. We deduce from formula (10) the relation

\[
\chi_ {a} (x) \equiv 1 + \theta (x) (a) \mathrm{T} \pmod {\mathrm{T} ^ {2} \mathrm{K} [ [ \mathrm{T} ] ]}\tag{11}
\]

for \(x\in \mathrm{R}_{\mathrm{K}}(\mathrm{A})\) and \(a\in \mathrm{A}\).

If E and \(E'\) are two A-modules of finite dimension over K with isomorphic semisimplifications, then we have \(\chi_{\mathrm{E}}(a;\mathrm{T})=\chi_{\mathrm{E}'}(a;\mathrm{T})\) .

THEOREM 2. --- Let \(\mathcal{A}\) be a generating subset of the K-vector space A. The homomorphism \(\chi_{\mathcal{A}}: \mathrm{R}_{\mathrm{K}}(\mathrm{A}) \to (1 + \mathrm{TK}[[\mathrm{T}]])^{\mathcal{A}}\) defined by the relation \(\chi_{\mathcal{A}}(x) = (\chi_a(x))_{a \in \mathcal{A}}\) is injective.

Let \(x\) be an element of \(\mathrm{R}_{\mathrm{K}}(\mathrm{A})\) such that \(\chi_{\mathcal{A}}(x) = 1\) . By (11), we have \(\theta(x)(a) = 0\) for every \(a \in \mathcal{A}\) , and therefore \(\theta(x) = 0\) because \(\theta(x)\) is a K-linear form on A and \(\mathcal{A}\) generates the K-vector space A. If the characteristic of K is zero, then this implies \(x = 0\) (VIII, p.~384, Corollary of Proposition 6), and therefore the result in this case. Suppose from now on that the characteristic \(p\) of K is nonzero.

Let us first treat the case when K is algebraically closed. By the above and loc. cit., we then have \(x \in p \mathrm{R}_{\mathrm{K}}(\mathrm{A})\) . Let \(y \in \mathrm{R}_{\mathrm{K}}(\mathrm{A})\) be such that x = py. For every element a of A, we have \(\chi_{a}(y)^{p} = \chi_{a}(py) = \chi_{a}(x) = 1\) . So we have \((\chi_{a}(y) - 1)^{p} = 0\) , and therefore \(\chi_{a}(y) = 1\) because the ring K{[}{[}T{]}{]} is an integral domain. So y belongs to the kernel of the endomorphism \(\chi_{A}\) . It follows by induction that x belongs to \(p^{n}\mathrm{R}_{\mathrm{K}}(\mathrm{A})\) for every integer \(n \geqslant 1\) . Since \(\mathrm{R}_{\mathrm{K}}(\mathrm{A})\) is a free Z-module, this implies x = 0, and therefore the injectivity of \(\chi_{A}\) in this case.

If K is no longer supposed to be algebraically closed, then we choose an algebraic closure \(\overline{K}\) of K and consider the diagram of groups and group homomorphisms

\[
\begin{array}{c} \mathrm{R} _ {\mathrm{K}} (\mathrm{A}) \xrightarrow {\chi_ {\mathcal {A}}} (1 + \mathrm{TK} [ [ \mathrm{T} ] ]) ^ {\mathcal {A}} \\ \Biggl \downarrow^ {u} \qquad \qquad \qquad \qquad \Biggl \downarrow^ {i} \\ \mathrm{R} _ {\overline {{\mathrm{K}}}} (\mathrm{A} _ {(\overline {{\mathrm{K}}})}) \xrightarrow {\overline {{\chi}} _ {\mathcal {A}}} (1 + \mathrm{T} \overline {{\mathrm{K}}} [ [ \mathrm{T} ] ]) ^ {\mathcal {A}}, \end{array}\tag{12}
\]

where \(u\) is the homomorphism deduced from the extension of scalars from K to \(\overline{\mathbf{K}}\) (VIII, p.~195), \(i\) is the canonical injection, and \(\overline{\chi}_{\mathcal{A}}\) is the homomorphism \(z \mapsto (\chi_{1 \otimes a}(z))_{a \in \mathcal{A}}\) . By formula (12) of III, §9, No.~1, p.~542, diagram (12) is commutative. By the above, the homomorphism \(\overline{\chi}_{\mathcal{A}}\) is injective. Since \(u\) is injective (VIII, p.~195, Theorem 1), the homomorphism \(\chi_{\mathcal{A}}\) is injective.

COROLLARY 1. --- Let E and F be semisimple A-modules of finite dimension over K, and let \(\mathcal{A}\) be a generating subset of the K-vector space A. Suppose that for every \(a \in \mathcal{A}\) , the characteristic polynomials of the endomorphisms \(a_{E}\) and \(a_{F}\) of the K-vector spaces E and F are equal. Then the A-modules E and F are isomorphic.

Let \(a\) be an element of \(\mathcal{A}\) . The characteristic polynomials of \(a_{\mathrm{E}}\) and \(a_{\mathrm{F}}\) have the same degree, so the dimension of E is equal to that of F; we denote it by \(n\) . Let \(\mathrm{Pc}_{\mathrm{E}}(a;\mathrm{T})\) be the characteristic polynomial of \(a_{\mathrm{E}}\) . In \(\mathrm{K(T)}\) , we have the equalities

\[
\chi_ {\mathrm{E}} (a; \mathrm{T}) = \det \left(1 + a _ {\mathrm{E}} \mathrm{T}\right) = (- \mathrm{T}) ^ {n} \det \left(\frac {- 1}{\mathrm{T}} - a _ {\mathrm{E}}\right) = (- \mathrm{T}) ^ {n} \operatorname{Pc} _ {\mathrm{E}} \left(a; \frac {- 1}{\mathrm{T}}\right),
\]

and \(\chi_{\mathrm{F}}(a;\mathrm{T})\) is given by an analogous formula. Because of our assumptions, we have \(\chi_{\mathrm{E}}(a;\mathrm{T})=\chi_{\mathrm{F}}(a;\mathrm{T})\) . By Theorem 2, we have \([E]=[F]\) , which implies that E and F are isomorphic (VIII, p.~190, Corollary of Proposition 7).

COROLLARY 2. --- Let A be a central simple algebra of finite degree over K. Let B be a semisimple K-algebra, let f and g be algebra homomorphisms from B to A, and let \(\mathcal{B}\) be a generating subset of the K-vector space B. The following properties are equivalent:

\begin{enumerate}
\def\labelenumi{(\roman{enumi})}
\item
  There exists an inner automorphism \(\theta\) of A such that \(g = \theta \circ f\) .
\item
  For every \(b \in \mathcal{B}\) , we have \(\mathrm{Pc}_{\mathrm{A / K}}(f(b); \mathrm{X}) = \mathrm{Pc}_{\mathrm{A / K}}(g(b); \mathrm{X})\) .
\end{enumerate}

When \(\mathrm{K}\) has characteristic zero, these properties are equivalent to the following:

\begin{enumerate}
\def\labelenumi{(\roman{enumi})}
\setcounter{enumi}{2}
\tightlist
\item
  For every \(b \in \mathcal{B}\) , we have \(\mathrm{Tr}_{\mathrm{A / K}}(f(b)) = \mathrm{Tr}_{\mathrm{A / K}}(g(b))\) .
\end{enumerate}

Denote by M and N the left B-modules obtained by endowing the additive group of A with the external laws \((b,a)\mapsto f(b)a\) and \((b,a)\mapsto g(b)a\) , respectively. By Proposition 1 of VIII, p.~257, property (i) is equivalent to the fact that the B-modules M and N are isomorphic. By construction, the homothety \(b_{M}\) associated with an element b of B is left multiplication by \(f(b)\) in A; consequently, we have the relations

\[
\begin{array}{c} \mathrm {Pc_ {M}} (b; \mathrm{X}) = \mathrm {Pc_ {A / K}} (f (b); \mathrm{X}), \\ \mathrm{Tr} _ {\mathrm{M}} (b) = \mathrm{Tr} _ {\mathrm{A/K}} (f (b)) \end{array}
\]

and two analogous relations where we replace M with N and f with g. We know (VIII, p.~386, Corollary 1) that the B-modules M and N are isomorphic if and only if we have \(\mathrm{Pc}_{\mathrm{M}}(b;\mathrm{X})=\mathrm{Pc}_{\mathrm{N}}(b;\mathrm{X})\) for every \(b\in B\) . When the field K has characteristic 0, this relation is also equivalent to \(\mathrm{Tr}_{\mathrm{M}}(b)=\mathrm{Tr}_{\mathrm{N}}(b)\) for every \(b\in B\) (VIII, p.~384, Corollary of Proposition 6). The equivalence of properties (i) and (ii), and also of (i) and (iii) when K has characteristic 0, follows.

\subsection{Coefficients of a Set of Classes of Modules}

Let C be a hereditary set of classes of A-modules (VIII, p.~183, Definition 1). We suppose that every A-module of type C is finite-dimensional over K and denote by \(\Theta_{\mathcal{C}}(\mathrm{A})\) the set of coefficients of the A-modules of type C. By Proposition 2 of VIII, p.~378, the set \(\Theta_{\mathcal{C}}(\mathrm{A})\) is also the union of the images of the A-linear mappings \(u: E \to A^{*}\) , where E runs through C; it is also the union of the subspaces \(\Theta_{\mathrm{E}}(\mathrm{A})\) of \(A^{*}\) , where E runs through C (loc. cit.). The family of (A, A)-sub-bimodules \((\Theta_{\mathrm{E}}(\mathrm{A}))_{\mathrm{E} \in \mathcal{C}}\) of \(A^{*}\) is directed, so its union \(\Theta_{\mathcal{C}}(\mathrm{A})\) is an (A, A)-sub-bimodule of \(A^{*}\) .

PROPOSITION 7. --- A left A-submodule of \(A^{*}\) of finite dimension over K is contained in \(\Theta_{\mathcal{C}}(\mathrm{A})\) if and only if it is of type C.

Let V be a left A-submodule of \(A^{*}\) of finite dimension over K. We saw in VIII, p.~379, Remark 3 that we have \(V \subset \Theta_{\mathrm{V}}(\mathrm{A})\) . If V is of type C, then we have \(\Theta_{\mathrm{V}}(\mathrm{A}) \subset \Theta_{\mathcal{C}}(\mathrm{A})\) , so V is contained in \(\Theta_{\mathcal{C}}(\mathrm{A})\) . Conversely, suppose that V is contained in \(\Theta_{\mathcal{C}}(\mathrm{A})\) . Since \(\Theta_{\mathcal{C}}(\mathrm{A})\) is the union of the directed family \((\Theta_{\mathrm{E}}(\mathrm{A}))_{\mathrm{E}\in\mathcal{C}}\) and V is finite-dimensional over K, there exists a module E of type C such that V is contained in \(\Theta_{\mathrm{E}}(\mathrm{A})\) . Now, there exists a natural number n such that \(\Theta_{\mathrm{E}}(\mathrm{A})\) is isomorphic to a quotient of \(\mathrm{E}^{n}\) (VIII, p.~378, Proposition 2). Since C is hereditary, V is of type C.

COROLLARY. --- Let E be an A-module of finite dimension over K. Then E is of type C if and only if \(\Theta_{\mathrm{E}}(\mathrm{A})\) is contained in \(\Theta_{\mathcal{C}}(\mathrm{A})\) .

The condition is obviously necessary.

Conversely, suppose that \(\Theta_{\mathcal{C}}(\mathrm{A})\) contains \(\Theta_{\mathrm{E}}(\mathrm{A})\) . The A-module \(\Theta_{\mathrm{E}}(\mathrm{A})\) is finite-dimensional over K (VIII, p.~378, Proposition 2), so it is of type C by Proposition 7. Now, there exists an integer \(n \geqslant 0\) such that E is isomorphic to an A-submodule of \(\Theta_{\mathrm{E}}(\mathrm{A})^{n}\) (VIII, p.~378, Proposition 2). Since C is hereditary, E is of type C.

\subsection{Cogebra Structure on the Restricted Dual}

For any \(a\) in \(\mathbf{A}\) , we denote by \(\eta(a)\) the linear form \(f \mapsto f(a)\) on \(\Theta(\mathbf{A})\) . We thus define a K-linear mapping \(\eta\) from \(\mathbf{A}\) to the dual \(\Theta(\mathbf{A})^*\) of the vector space \(\Theta(\mathbf{A})\) . Set \(\varepsilon = \eta(1)\) .

PROPOSITION 8. --- On the vector space \(\Theta(A)\) , there exists a unique cogebra structure (III, §11, No.~1, p.~574) such that the mapping \(\eta: A \to \Theta(A)^*\) is a homomorphism from \(A\) to the dual algebra (III, §11, No.~2, p.~579) of \(\Theta(A)\) . The cogebra \(\Theta(A)\) is coassociative and admits \(\varepsilon\) as a counit.

For every integer \(n \geqslant 1\) , consider the K-linear mapping \(j_{n}\) from \(\Theta(\mathrm{A})^{\otimes n}\) to the dual of \(A^{\otimes n}\) characterized by the formula

\[
\langle j _ {n} \left(f _ {1} \otimes \dots \otimes f _ {n}\right), a _ {1} \otimes \dots \otimes a _ {n} \rangle = \prod_ {i = 1} ^ {n} \left\langle f _ {i}, a _ {i} \right\rangle\tag{13}
\]

for \((a_{i})\) in \(\mathrm{A}^n\) and \((f_i)\) in \(\Theta (\mathrm{A})^n\) . By Proposition 16, (ii) of II, §7, No.~7, p.~308, the mapping \(j_{n}\) is injective. Denote by \(m_{\mathrm{K}}:\mathrm{K}\otimes \mathrm{K}\to \mathrm{K}\) and \(m_{\mathrm{A}}:\mathrm{A}\otimes \mathrm{A}\rightarrow \mathrm{A}\) the mappings deduced from the multiplication in K and A, respectively. For \(f,g\in \Theta (\mathrm{A})\) and \(a,b\in \mathrm{A}\) , we have

\[
\langle j _ {2} (f \otimes g), a \otimes b \rangle = m _ {\mathrm{K}} \circ (\eta (a) \otimes \eta (b)) (f \otimes g).
\]

We therefore have

\[
\langle j _ {2} (t), a \otimes b \rangle = m _ {\mathrm{K}} \circ (\eta (a) \otimes \eta (b)) (t)\tag{14}
\]

for every \(t\in \Theta (\mathrm{A})\otimes \Theta (\mathrm{A})\).

Lemma 2. --- Let \(c: \Theta(A) \to \Theta(A) \otimes \Theta(A)\) be a K-linear mapping. Then \(\eta\) is a homomorphism from A to the dual algebra of the cogebra \((\Theta(A), c)\) if and only if the following diagram commutes:

\[
\begin{array}{c} \Theta (\mathrm{A}) \xrightarrow {c} \Theta (\mathrm{A}) \otimes \Theta (\mathrm{A}) \\ \Biggl \downarrow_ {j _ {1}} \qquad \qquad \qquad \qquad \Biggl \downarrow_ {j _ {2}} \\ \mathrm{A} ^ {*} \xrightarrow {^ t m _ {\mathrm{A}}} (\mathrm{A} \otimes \mathrm{A}) ^ {*}. \end{array}\tag{15}
\]

Indeed, \(\eta\) is a homomorphism from A to the dual algebra of the cogebra \((\Theta (\mathrm{A}),c)\) if and only if we have \(\eta (ab) = m_{\mathrm{K}}\circ (\eta (a)\otimes \eta (b))\circ c\) for all \(a,b\in \mathrm{A}\) , that is

\[
\eta (a b) (f) = m _ {\mathrm{K}} \circ (\eta (a) \otimes \eta (b)) (c (f))
\]

for \(a, b \in \mathrm{A}\) and \(f \in \Theta(\mathrm{A})\) . Now, we have

\[
\eta (a b) (f) = f (a b) = \left\langle {} ^ {t} m _ {\mathrm{A}} \left(j _ {1} (f)\right), a \otimes b \right\rangle
\]

and, by (14),

\[
m _ {K} \circ (\eta (a) \otimes \eta (b)) (c (f)) = \langle j _ {2} (c (f)), a \otimes b \rangle
\]

for all \(a, b \in \mathrm{A}\) and every \(f \in \Theta(\mathrm{A})\) . The lemma follows.

Since \(j_{2}\) is injective, there exists at most one linear mapping c that makes the diagram above commute. To prove its existence, we must prove that the image of \(^{t}m \circ j_{1}\) is contained in that of \(j_{2}\) . In other words, we must prove that there exist, for every element f of \(\Theta(\mathrm{A})\) , a natural number n and elements \(f_{1}^{\prime}, \ldots, f_{n}^{\prime}, f_{1}^{\prime\prime}, \ldots, f_{n}^{\prime\prime}\) of \(\Theta(\mathrm{A})\) satisfying the relations

\[
f (a b) = \sum_ {i = 1} ^ {n} f _ {i} ^ {\prime} (a) f _ {i} ^ {\prime \prime} (b)\tag{16}
\]

for \(a, b \in \mathbf{A}\) . We will then have

\[
c (f) = \sum_ {i = 1} ^ {n} f _ {i} ^ {\prime} \otimes f _ {i} ^ {\prime \prime}.\tag{17}
\]

By the corollary of VIII, p.~379, there exists a left A-module E of finite dimension over K with f as a coefficient. Let \((e_{1},\ldots,e_{n})\) be a basis of E, \((e_{1}^{*},\ldots,e_{n}^{*})\) the dual basis, x an element of E, and \(x^{*}\) an element of \(E^{*}\) such that \(f=c_{\mathrm{E}}(x,x^{*})\) . Set \(f_{i}^{\prime}=c_{\mathrm{E}}(e_{i},x^{*})\) and \(f_{i}^{\prime\prime}=c_{\mathrm{E}}(x,e_{i}^{*})\) for \(i\in[1,n]\) ; for

\(a, b\) in A, we have

\[
\begin{array}{r l} & {f (a b) = \langle x ^ {*}, a b x \rangle = \langle x ^ {*} a, b x \rangle = \left\langle \sum_ {i} \langle x ^ {*} a, e _ {i} \rangle e _ {i} ^ {*}, b x \right\rangle} \\ & {\qquad = \sum_ {i} \langle x ^ {*} a, e _ {i} \rangle \langle e _ {i} ^ {*}, b x \rangle = \sum_ {i} \langle x ^ {*}, a e _ {i} \rangle \langle e _ {i} ^ {*}, b x \rangle = \sum_ {i} f _ {i} ^ {\prime} (a) f _ {i} ^ {\prime \prime} (b),} \end{array}
\]

and therefore (16).

Let us prove the coassociativity of \(c\) . For this, consider the K-linear mappings

\[
c ^ {\prime} = \left(c \otimes 1 _ {\Theta (\mathrm{A})}\right) \circ c \quad \text { and } \quad c ^ {\prime \prime} = \left(1 _ {\Theta (\mathrm{A})} \otimes c\right) \circ c
\]

from \(\Theta (\mathrm{A})\) to \(\Theta (\mathrm{A})^{\otimes 3}\) . We have the relations

\[
\begin{array}{r l} \langle j _ {3} (f \otimes c (g)), a \otimes b \otimes c \rangle & = \langle f, a \rangle \langle j _ {2} \circ c (g), b \otimes c \rangle \\ & = \langle f, a \rangle \langle g, b c \rangle \\ & = \langle j _ {2} (f \otimes g), a \otimes b c \rangle \\ & = \langle^ {t} (\mathrm{Id} _ {\mathrm{A}} \otimes m _ {\mathrm{A}}) \circ j _ {2} (f \otimes g), a \otimes b \otimes c \rangle \end{array}
\]

for \(f, g \in \Theta(\mathbf{A})\) and \(a, b, c \in \mathbf{A}\) . From this, we deduce that the following diagram commutes:

\[
\begin{array}{c} \Theta (\mathrm{A}) \otimes \Theta (\mathrm{A}) \xrightarrow {\operatorname{Id} _ {\Theta (\mathrm{A})} \otimes c} \Theta (\mathrm{A}) \otimes \Theta (\mathrm{A}) \otimes \Theta (\mathrm{A}) \\ \Biggl \downarrow_ {j _ {2}} \\ (\mathrm{A} \otimes \mathrm{A}) ^ {*} \xrightarrow {^ t (\operatorname{Id} _ {\mathrm{A}} \otimes m _ {\mathrm{A}})} (\mathrm{A} \otimes \mathrm{A} \otimes \mathrm{A}) ^ {*}. \end{array}\tag{18}
\]

Because of the commutativity of this diagram and that of (15), for \(f \in \Theta(\mathrm{A})\) and \(a, a', a'' \in A\) , we have

\[
\langle j _ {3} \circ c ^ {\prime} (f), a \otimes a ^ {\prime} \otimes a ^ {\prime \prime} \rangle = \langle f, (a a ^ {\prime}) a ^ {\prime \prime} \rangle ;
\]

we can show the relation

\[
\langle j _ {3} \circ c ^ {\prime \prime} (f), a \otimes a ^ {\prime} \otimes a ^ {\prime \prime} \rangle = \langle f, a (a ^ {\prime} a ^ {\prime \prime}) \rangle
\]

likewise. Since multiplication in A is associative, we have \(j_{3} \circ c' = j_{3} \circ c''\) , and therefore \(c' = c''\) because \(j_{3}\) is injective.

Finally, formulas (16) and (17) imply that \(\Theta(\mathrm{A})\) admits \(\varepsilon\) as a counit.

Remark 1. --- Let \((\mathrm{V},\pi)\) be a finite-dimensional linear representation of the algebra A. Let us introduce a basis \((e_{1},\ldots,e_{n})\) of V and the dual basis \((e_{1}^{*},\ldots,e_{n}^{*})\) of \(V^{*}\) . By the proof of Lemma 2, we have the relation

\[
c (c _ {\pi} (x, x ^ {*})) = \sum_ {k = 1} ^ {n} c _ {\pi} (e _ {k}, x ^ {*}) \otimes c _ {\pi} (x, e _ {k} ^ {*})\tag{19}
\]

for \(x \in V\) and \(x^{*} \in V^{*}\) . For \(1 \leqslant i, j \leqslant n\) , set \(\pi_{ij} = c_{\pi}(e_j, e_i^*)\) . For every \(a \in A\) , the matrix of \(\pi(a)\) with respect to the basis \((e_1, \ldots, e_n)\) of V is equal to \((\pi_{ij}(a))\) . For \(1 \leqslant i \leqslant n\) and \(1 \leqslant j \leqslant n\) , we then have

\[
c (\pi_ {i j}) = \sum_ {k = 1} ^ {n} \pi_ {i k} \otimes \pi_ {k j}.\tag{20}
\]

DEFINITION 3. --- Let C be a cogebra over the field K and c its coproduct. A subcogebra of C is any linear subspace \(C_{1}\) of C such that \(c(\mathrm{C}_{1})\) is contained in the canonical image of \(C_{1} \otimes_{K} C_{1}\) in \(C \otimes_{K} C\) .

Let \(j\) be the canonical injection of \(C_1\) into \(C\) . By this definition, there exists a unique linear mapping \(c_1: C_1 \to C_1 \otimes_K C_1\) such that we have

\[
c \circ j = (j \otimes j) \circ c _ {1};\tag{21}
\]

this relation means that \(j\) is a morphism of cogebras from \((\mathrm{C}_1, c_1)\) to \((\mathrm{C}, c)\) (III, §11, No.~1, p.~574).

If C is coassociative, then \(C_{1}\) is coassociative. If C is cocommutative, then so is \(C_{1}\) . If C has a counit \(\varepsilon\) , then the restriction of \(\varepsilon\) to \(C_{1}\) is a counit of \(C_{1}\) .

PROPOSITION 9. --- Let \(\Theta\) be a linear subspace of \(\Theta(A)\) . The following properties are equivalent:

\begin{enumerate}
\def\labelenumi{(\roman{enumi})}
\item
  \(\Theta\) is an (A, A)-sub-bimodule of \(\Theta(A)\) .
\item
  \(\Theta\) is a subcogebra of \(\Theta(A)\) .
\item
  There exists a hereditary set \(\mathcal{C}\) of classes of A-modules of finite dimension over \(K\) such that \(\Theta = \Theta_{\mathcal{C}}(A)\) .
\end{enumerate}

When these properties hold, the set C mentioned in (iii) is uniquely determined.

The last assertion follows from the corollary of VIII, p.~388: the set C consists of the classes of A-modules E of finite dimension over K such that \(\Theta_{\mathrm{E}}(\mathrm{A})\) is contained in \(\Theta\) .

\begin{enumerate}
\def\labelenumi{(\roman{enumi})}
\setcounter{enumi}{2}
\item
  \(\Rightarrow\) (ii): Let C be a hereditary set of classes of A-modules of finite dimension over K. Then \(\Theta_{\mathcal{C}}(\mathrm{A})\) is the union of the directed family \((\Theta_{\mathrm{E}}(\mathrm{A}))_{\mathrm{E}\in\mathcal{C}}\) . Since \(\Theta_{\mathrm{E}}(\mathrm{A})\) is a subcogebra of \(\Theta(\mathrm{A})\) for every \(E\in C\) (VIII, p.~390, formula (19)), the same holds for \(\Theta_{\mathcal{C}}(\mathrm{A})\) .
\item
  \(\Rightarrow\) (i): Let \(f\in \Theta (\mathrm{A})\) . Let \(f_{1}^{\prime},\ldots ,f_{n}^{\prime},f_{1}^{\prime \prime},\ldots ,f_{n}^{\prime \prime}\) be elements of \(\Theta (\mathrm{A})\) satisfying \(c(f) = \sum f_i^\prime \otimes f_i''\) . For \(a,b\) in A, we have \(f(ab) = \sum f_i'(a)f_i''(b)\) , and therefore
\end{enumerate}

\[
b f = \sum_ {i = 1} ^ {n} f _ {i} ^ {\prime \prime} (b) f _ {i} ^ {\prime}, \quad f a = \sum_ {i = 1} ^ {n} f _ {i} ^ {\prime} (a) f _ {i} ^ {\prime \prime}
\]

(VIII, p.~375, formulas (3) and (4)). Consequently, a subcogebra of \(\Theta(A)\) is an \((A, A)\) -sub-bimodule.

\begin{enumerate}
\def\labelenumi{(\roman{enumi})}
\tightlist
\item
  \(\Rightarrow\) (iii): Suppose that \(\Theta\) is an (A, A)-sub-bimodule of \(\Theta(\mathrm{A})\) ; let C be the set of classes of A-modules E of finite dimension over K such that \(\Theta_{\mathrm{E}}(\mathrm{A})\) is contained in \(\Theta\) . The set C is hereditary (VIII, p.~377, Remark 2), and we have \(\Theta_{\mathcal{C}}(\mathrm{A}) \subset \Theta\) by construction. Let \(f \in \Theta\) and E = Af. Then E is finite-dimensional over K. Consequently, every linear form on E is of the form \(u_{a}: g \mapsto g(a)\) with a in A (II, §7, No.~5, p.~302, Corollary 2). Now, for \(a, b, x \in A\) , we have
\end{enumerate}

\[
c _ {\mathrm{E}} (a f, u _ {b}) (x) = \langle u _ {b}, x a f \rangle = f (b x a) = a f b (x),
\]

so that \(c_{\mathrm{E}}(af, u_{b}) = afb\) . We therefore have \(\Theta_{\mathrm{E}}(\mathrm{A}) \subset \Theta\) . Consequently, the A-module E is of type C, and f is one of its coefficients. We therefore have \(\Theta \subset \Theta_{\mathcal{C}}(\mathrm{A})\) and, finally, \(\Theta = \Theta_{\mathcal{C}}(\mathrm{A})\) .

Remark 2. --- Let \(\Theta\) be a subcogebra of \(\Theta(A)\) . We endow \(K\) with the discrete topology and the algebra \(A\) with the coarsest topology for which the mappings \(f: A \to K\) for \(f\) running through \(\Theta\) are continuous (Gen.~Top., I, §2, No.~2, p.~175). This topology endows \(A\) with the structure of a topological \(K\) -module. A two-sided ideal \(\mathfrak{a}\) of \(A\) is open if and only if it has finite codimension and its orthogonal \(\mathfrak{a}'\) is contained in \(\Theta\) . By Proposition 3 of VIII, p.~379, the open two-sided ideals of \(A\) form a fundamental system of neighborhoods of \(0\) in \(A\) . The topology on \(A\) is therefore compatible with its ring structure.

Let C be the hereditary set of classes of A-modules of finite dimension over K such that \(\Theta = \Theta_{C}\) (VIII, p.~391, Proposition 9). Let E be a left A-module of finite dimension over K. We endow it with the discrete topology. The following properties are equivalent:

\begin{enumerate}
\def\labelenumi{(\roman{enumi})}
\item
  The A-module E is of type C.
\item
  The annihilator of the A-module E is open in A.
\item
  The mapping \((a, x) \mapsto ax\) from A × E to E is continuous.
\end{enumerate}

The last property means that \(\mathrm{E}\) is a topological A-module.

Let \(\Theta^{*}\) be the dual algebra of the cogebra \(\Theta\) . We endow \(\Theta^{*}\) with the coarsest topology for which the mappings \(\varphi \mapsto \varphi(u)\) from \(\Theta^{*}\) to K, for \(u\) running through \(\Theta\) , are continuous. The topology on the algebra \(\Theta^{*}\) is compatible with the additive group structure on \(\Theta^{*}\) . The orthogonals in \(\Theta^{*}\) of the sets of the form \(\Theta_{\mathrm{E}}(\mathrm{A})\) , where E is an A-module of type \(\mathcal{C}\) , form a fundamental system of neighborhoods of 0. Now, such a set is a subcogebra of \(\Theta\) , so its orthogonal is an ideal of \(\Theta^{*}\) . The topology on \(\Theta^{*}\) is therefore compatible with its ring structure. The canonical algebra homomorphism \(\eta: A \to \Theta^{*}\) (that sends \(a \in A\) to the linear form \(f \mapsto f(a)\) on \(\Theta\) ) defines an isomorphism from the completed Hausdorff space \(\widehat{A}\) of A (Gen.~Top., II, §3, No.~7, p.~191) to \(\Theta^{*}\) .

\BourbakiExercises{Exercises}

\begin{enumerate}
\def\labelenumi{\arabic{enumi})}
\tightlist
\item
  Let \(\mathrm{K}\) be a commutative field. Let \(\mathrm{A}\) be the K-algebra with basis consisting of the unit element and two elements \(e, f\) such that \(e^2 = e\) , \(ef = f\) , and \(fe = f^2 = 0\) . a) Prove that the regular and coregular representations of \(\mathrm{A}\) are not isomorphic (calculate \(\operatorname{Tr}_{\mathrm{A / K}}(e)\) and \(\operatorname{Tr}_{\mathrm{A}^{\circ} / \mathrm{K}}(e)\) ).
\end{enumerate}

\begin{enumerate}
\def\labelenumi{\alph{enumi})}
\setcounter{enumi}{1}
\tightlist
\item
  Prove that the regular representation contains a subrepresentation that is isomorphic to neither a subrepresentation nor a quotient representation of the coregular representation. Moreover, there exists a simple representation of A that is not isomorphic to a subrepresentation of the regular representation.
\end{enumerate}

\begin{enumerate}
\def\labelenumi{\arabic{enumi})}
\setcounter{enumi}{1}
\tightlist
\item
  Let \(A\) be a K-algebra. We view the dual \(A^*\) of \(A\) as a left \(A\) -module.
\end{enumerate}

\begin{enumerate}
\def\labelenumi{\alph{enumi})}
\item
  The mapping \(E \mapsto E'\) is a surjection from the set of A-submodules of \(A^{*}\) to the set of right ideals of A. When E is such a submodule, the orthogonal of \(E'\) in \(A^{*}\) is a submodule of \(A^{*}\) containing E.
\item
  Prove that the orthogonal in A of the socle of \(\mathbf{A}^*\) is the radical of A.
\item
  Prove that the orthogonal in \(A^{*}\) of the socle of \(A_{d}\) is the radical of the A-module \(A^{*}\) . Deduce that \(A^{*}\) is a semisimple A-module if and only if the ring A is semisimple.
\end{enumerate}

¶ 3) Let A be an Artinian ring and \(\overline{\mathrm{A}}\) the ring \(\mathrm{A} / \Re (\mathrm{A})\) . We take a decomposition of A as a direct sum of indecomposable ideals \(\mathrm{Ae}_i\) ( \(i\in \mathrm{I}\) ), where the \(e_i\) are orthogonal idempotents (VIII, p.~179, Exercise 6, d)). We denote by \(\mathrm{I} = \cup_{\alpha \in \mathrm{R}}\mathrm{I}_{\alpha}\) the partition of I such that two elements \(i,j\) of I belong to the same set \(\mathrm{I}_{\alpha}\) if and only if the modules \(\mathrm{Ae}_i\) and \(\mathrm{Ae}_j\) are isomorphic. The \(\overline{\mathrm{A}}\) -modules \(\bigoplus_{i\in \mathrm{I}_{\alpha}}\overline{\mathrm{A}}\overline{e}_i\) are the isotypical components of \(\overline{\mathrm{A}}\) (VIII, p.~179, Exercises 6, d) and 7, a))

\begin{enumerate}
\def\labelenumi{\alph{enumi})}
\item
  Prove that if the ring A is autoinjective (VIII, p.~53, Exercise 8), then there exists a permutation \(\sigma\) of R such that every right ideal \(e_i\mathrm{A}\) , for \(i\in \mathrm{I}_{\alpha}\) , contains a unique minimal right ideal isomorphic to \(\overline{e}_j\overline{\mathrm{A}}\) for \(j\in \mathrm{I}_{\sigma (\alpha)}\) (observe that the right annihilator of an idempotent \(e\) is \((1 - e)\mathrm{A}\) ). Deduce that the left socle \(\mathfrak{s}\) of A is contained in its right socle \(\mathfrak{t}\) (prove that every right ideal \(e_i\mathfrak{s}\) is minimal, by proving that it is the annihilator of a maximal left ideal, cf.~VIII, p.~178, Exercise 2). Conclude that we have \(\mathfrak{s} = \mathfrak{t}\) and that every left ideal \(\mathrm{Ae}_i\) , for \(i\in \mathrm{I}_{\sigma (\alpha)}\) , contains a unique minimal left ideal isomorphic to \(\overline{\mathrm{A}}\overline{e}_j\) for every \(j\in \mathrm{I}_{\alpha}\) .
\item
  Conversely, if every left ideal \(Ae_{i}\) contains a unique minimal left ideal and every right ideal \(e_{i}A\) contains a unique minimal right ideal, then the ring A is autoinjective. (Use Exercise 7 of VIII, p.~179 to prove that every foot of s is a minimal two-sided ideal. Deduce that t contains s, and then that s = t. Finally, prove that condition \(\beta\) ) of Exercise 8 of VIII, p.~53 is satisfied.)
\item
  Suppose that A is an algebra of finite degree over a commutative field and an autoinjective ring. Let \(\alpha \in R\) , \(i \in I_{\alpha}\) , and \(j \in I_{\sigma(\alpha)}\) . Prove that the linear representation of A associated with the A-module \(Ae_{j}\) is isomorphic to the transpose of the linear representation of \(A^{\circ}\) associated with the right A-module \(e_{i}A\) . (Let \(U_{i}\) be the left submodule of \(A^{*}\) that is the orthogonal of \((1 - e_{i})A\) . Deduce from Exercise 7, d) of VIII, p.~179 that \(U_{i}\) is isomorphic to a quotient module of \(Ae_{j}\) and that we consequently have long \(e_{i}A \leqslant long e_{j}A\) . Conclude that these lengths are equal, so \(U_{i}\) is isomorphic to \(Ae_{j}\) . Use b) and Exercise 2, a) to prove a converse of this statement.
\end{enumerate}

\begin{enumerate}
\def\labelenumi{\arabic{enumi})}
\setcounter{enumi}{3}
\tightlist
\item
  Keep the assumptions of Exercise 3.
\end{enumerate}

\begin{enumerate}
\def\labelenumi{\alph{enumi})}
\item
  Prove that if A is a Frobenius ring (VIII, p.~53, Exercise 9), then we have \(\operatorname{Card}(\mathrm{I}_{\alpha}) = \operatorname{Card}(\mathrm{I}_{\sigma(\alpha)})\) for every \(\alpha \in R\) . Also prove the converse.
\item
  An algebra of finite degree over a commutative field is Frobenius if and only if its regular representation is isomorphic to its coregular representation (use Exercise 3, c)).
\end{enumerate}

\begin{enumerate}
\def\labelenumi{\arabic{enumi})}
\setcounter{enumi}{4}
\tightlist
\item
  Let K be a commutative field and A be a K-algebra. We endow \(A^{*}\) with its canonical left A-module structure. For \(t \in A^{*}\) , we denote by \(\varphi_{t}\) the mapping \((x, y) \mapsto t(xy)\) from \(A \times A\) to K.
\end{enumerate}

\begin{enumerate}
\def\labelenumi{\alph{enumi})}
\item
  Prove that this defines a K-linear isomorphism from \(A^{*}\) to the set of K-bilinear mappings \(\varphi: A \times A \to K\) satisfying \(\varphi(xy, z) = \varphi(x, yz)\) for \(x, y, z \in A\) .
\item
  We say that the element \(t \in A^{*}\) is invertible if the homomorphism \(a \mapsto at\) from A to \(A^{*}\) is bijective. Let A be a K-algebra of finite degree. Prove that the following conditions are equivalent:
\end{enumerate}

\begin{enumerate}
\def\labelenumi{(\roman{enumi})}
\item
  A is Frobenius (Exercise 4).
\item
  A admits an invertible linear form.
\item
  A admits a linear form whose kernel does not contain any left ideals of \(\mathbf{A}\) .
\end{enumerate}

\begin{enumerate}
\def\labelenumi{\arabic{enumi})}
\setcounter{enumi}{5}
\tightlist
\item
  Let K be a commutative field. We say that a K-algebra A is symmetric if it admits a tracial invertible linear form (VIII, p.~115, Exercise 3).
\end{enumerate}

\begin{enumerate}
\def\labelenumi{\alph{enumi})}
\item
  Prove that a symmetric algebra has finite degree and is Frobenius.
\item
  Prove that the algebra of a finite group over \(\mathrm{K}\) is symmetric.
\item
  A semisimple algebra of finite degree over \(\mathrm{K}\) is symmetric.
\end{enumerate}

\section{Linear Representations of Finite Groups}

In this section, G is a group, and K is a commutative ring. If G is a finite group, then we denote by \(|G|\) the element \(\operatorname{Card}(G) \cdot 1\) of K. From No.~5 on, we assume that the group G is finite, that the ring K is an algebraically closed field, and that \(|G|\) is not zero.

\subsection{Linear Representations}

DEFINITION 1. --- Let M be a K-module. A linear representation of G in M is a group homomorphism from G to the linear group GL(M) (II, §1, No.~2, p.~195).

Let \(\pi\) be a linear representation of G in a K-module M. We also say that the pair \((M,\pi)\) is a linear representation of G. When K is nonzero and M is a free K-module, the dimension of M is called the degree (or dimension) of the representation \(\pi\) .

Recall that the canonical mapping \(g \mapsto e_{g}\) from G to the algebra K{[}G{]} of the group G is a basis of the K-module K{[}G{]}. By abuse of notation, we also denote the element \(e_{g}\) of K{[}G{]} by g. Given a linear representation \(\pi : G \to \mathbf{GL}(M)\) of G, there exists a unique homomorphism of K-algebras from K{[}G{]} to \(\operatorname{End}_{\mathrm{K}}(\mathrm{M})\) that extends \(\pi\) (III, §2, No.~6, p.~447, Example); we also denote this extension by \(\pi\) . It is a representation of the algebra K{[}G{]} in M (VIII, p.~373, Definition 1), so that M is a K{[}G{]}-module. Conversely, every linear representation \(\rho\) of the algebra K{[}G{]} in M defines a linear representation of G in M given by \(g \mapsto \rho(g)\) .

We will freely apply the terminology concerning the K{[}G{]}-module structure to linear representations of the group G. For example, we speak of a subrepresentation, a simple representation, a direct sum of representations, etc. Let \((\mathrm{M},\pi)\) and \((\mathrm{M}^{\prime},\pi^{\prime})\) be linear representations of G; we denote by \(\operatorname{Hom}_{\mathrm{G}}(\pi,\pi^{\prime})\) the K-module \(\operatorname{Hom}_{\mathrm{K[G]}}(\mathrm{M},\mathrm{M}^{\prime})\) of homomorphisms of K{[}G{]}-modules from M to M'. \(^{(1)}\)

A mapping \(f\) from \(G\) to \(K\) is called a central function if we have \(f(gg') = f(g'g)\) for every pair \((g, g')\) of elements of \(G\) ; equivalently, we have \(f(ghg^{-1}) = f(h)\) for every pair of elements \(g, h\) of \(G\) . The central functions are therefore the functions on \(G\) whose restriction to each conjugacy class is constant. They form a submodule of the space of mappings from \(G\) to \(K\) , denoted by \(\mathcal{Z}_K(G)\) . When \(G\) is finite, the K-algebra \(\mathcal{Z}_K(G)\) is a free K-module of dimension the number of conjugacy classes of \(G\) . The center \(Z(K[G])\) of the algebra \(K[G]\) consists of the elements \(a = \sum a_g g\) of \(K[G]\) such that \(hah^{-1} = a\) for every \(h \in G\) . Now, we have

\[
h a h ^ {- 1} = \sum_ {g \in \mathrm{G}} a _ {h ^ {- 1} g h} g;
\]

consequently, when G is finite, the center \(Z(K[G])\) of the algebra K{[}G{]} consists of the central functions.

Let \((\mathrm{M},\pi)\) be a linear representation of G. Suppose that M is a finite-dimensional free K-module. The trace of \(\pi\) is the trace of the representation of K{[}G{]} associated with \(\pi\) , that is (VIII, p.~382), the linear form \(a\mapsto\mathrm{Tr}(\pi(a))\) on K{[}G{]}. This linear form is determined by the mapping \(g\mapsto\mathrm{Tr}(\pi(g))\) from G to K, which we call the character of the representation \(\pi\) and denote by \(\chi_{\pi}\) . The character of a representation is a central function (II, §4, No.~3, p.~273, Proposition 3).

Let \(M\) and \(M'\) be finite-dimensional free \(K\) -modules. Let \(\pi\) and \(\pi'\) be linear representations of \(G\) in \(M\) and \(M'\) , respectively. Then \(M \oplus M'\) is a finite-dimensional free \(K\) -module, and, by Proposition 1 of III, §9, No.~2, p.~543, we have

\[
\chi_ {\pi \oplus \pi^ {\prime}} = \chi_ {\pi} + \chi_ {\pi^ {\prime}}.
\]

More generally, let

\[
0 \longrightarrow \mathrm{M} ^ {\prime} \longrightarrow \mathrm{M} \longrightarrow \mathrm{M} ^ {\prime \prime} \longrightarrow 0
\]

be an exact sequence of K{[}G{]} modules, and let \(\pi\) , \(\pi'\) , and \(\pi''\) be the linear representations of G associated with M, \(M'\) , and \(M''\) , respectively. Suppose that the K-modules \(M'\) and \(M''\) are free and finite-dimensional. Then so is M (II, §1, No.~11, p.~218, Proposition 21), and we have

\[
\chi_ {\pi} = \chi_ {\pi^ {\prime}} + \chi_ {\pi^ {\prime \prime}}.\tag{1}
\]

PROPOSITION 1. --- Suppose that K is a commutative field. Let \(\pi\) and \(\pi'\) be finite-dimensional semisimple K-linear representations of G.

\begin{enumerate}
\def\labelenumi{\alph{enumi})}
\item
  If the characteristic polynomials of \(\pi(g)\) and \(\pi'(g)\) are the same for every \(g \in G\) , then the representations \(\pi\) and \(\pi'\) are isomorphic.
\item
  Suppose, moreover, that K has characteristic 0. If the characters \(\chi_{\pi}\) and \(\chi_{\pi'}\) are equal, then the representations \(\pi\) and \(\pi'\) are isomorphic.
\end{enumerate}

Assertion a) follows from Corollary 1 of VIII, p.~386; assertion b) follows from part a) of the corollary of VIII, p.~384.

Examples. --- 1) The unit or trivial representation of G is the representation \((\mathrm{K}, \varepsilon)\) , where \(\varepsilon(g) = \operatorname{Id}_{\mathrm{K}}\) for every \(g \in G\) . Its character is the constant function with value 1.

\begin{enumerate}
\def\labelenumi{\arabic{enumi})}
\setcounter{enumi}{1}
\tightlist
\item
  The (left) regular representation of G is the representation \(\gamma\) of G in K{[}G{]} defined by \(\gamma(g)(x)=gx\) for \(g\in G\) and \(x\in K[G]\) . It corresponds to the left regular representation of the algebra K{[}G{]} (VIII, p.~375). Suppose that the group G is finite. The regular representation then has finite degree. For every element g of G different from the identity, the matrix of left multiplication by g in K{[}G{]} with respect to the canonical basis is the matrix of a permutation without fixed point. We therefore have
\end{enumerate}

\[
\chi_ {\boldsymbol {\gamma}} (g) = \left\{ \begin{array}{l l} | \mathrm{G} | & \text { if   g   is   the   identity   element }, \\ 0 & \text { otherwise }. \end{array} \right.\tag{2}
\]

We define the right regular representation of G likewise. The biregular representation of G is the representation \(\rho\) of G × G in K{[}G{]} defined by \(\rho(g, g')(x) = gxg'^{-1}\) for \((g, g') \in \mathrm{G} \times \mathrm{G}\) and \(x \in \mathrm{K}[\mathrm{G}]\) .

\begin{enumerate}
\def\labelenumi{\arabic{enumi})}
\setcounter{enumi}{2}
\item
  Given a linear representation \((\mathrm{M},\pi)\) of G, the contragredient or dual representation \(\pi^{\vee}\) of \(\pi\) is the representation of G in the K-module dual to M defined by the relation \(\pi^{\vee}(g)={}^{t}\pi(g^{-1})\) for every \(g\in G\) (cf.~II, §2, No.~5, p.~234). If M is a finite-dimensional free K-module, then so is its dual, and we have \(\chi_{\pi^{\vee}}(g)=\chi_{\pi}(g^{-1})\) for every \(g\in G\) .
\item
  Let \((\mathrm{M},\pi)\) and \((\mathrm{M}^{\prime},\pi^{\prime})\) be linear representations of G. In Example 1 of VIII, p.~198, we defined a K{[}G{]}-module structure on \(M\otimes_{K}M^{\prime}\) . The corresponding linear representation is called the tensor product of \(\pi\) and \(\pi^{\prime}\) and is denoted by \(\pi\otimes\pi^{\prime}\) . For \(g\in G\) , \(x\in M\) , and \(x^{\prime}\in M^{\prime}\) , we have \(\big((\pi\otimes\pi^{\prime})(g)\big)(x\otimes x^{\prime})=\pi(g)x\otimes\pi^{\prime}(g)x^{\prime}\) .
\end{enumerate}

If M and M' are finite-dimensional free K-modules, then by Proposition 2 of III, §9, No.~2, p.~543, we have

\[
\chi_ {\pi \otimes \pi^ {\prime}} = \chi_ {\pi} \chi_ {\pi^ {\prime}}.\tag{3}
\]

\begin{enumerate}
\def\labelenumi{\arabic{enumi})}
\setcounter{enumi}{4}
\tightlist
\item
  Suppose that G is the product \(G' \times G''\) of two groups. The K-linear mapping from \(K[G'] \otimes_{K} K[G'']\) to K{[}G{]} that sends \(g' \otimes g''\) to \((g', g'')\) for \(g' \in G'\) and \(g'' \in G''\) is an algebra isomorphism. Let \((M', \pi')\) be a linear representation of \(G'\) and \((M'', \pi'')\) a linear representation of \(G''\) . The external tensor product of \(\pi'\) and \(\pi''\) , denoted by \(\pi' \boxtimes \pi''\) , is the representation of \(G' \times G''\) in the vector space \(M' \otimes M''\) defined by \((\pi' \boxtimes \pi'')(g', g'') = \pi'(g') \otimes \pi''(g'')\) for \((g', g'') \in G' \times G''\) . If \(M'\) and \(M''\) are finite-dimensional free K-modules, then \(M' \otimes_{K} M''\) is a finite-dimensional free K-module and the character of the representation \(\pi' \boxtimes \pi''\) is given by the formula
\end{enumerate}

\[
\chi_ {\pi^ {\prime} \boxtimes \pi^ {\prime \prime}} (g ^ {\prime}, g ^ {\prime \prime}) = \chi_ {\pi^ {\prime}} (g ^ {\prime}) \chi_ {\pi^ {\prime \prime}} (g ^ {\prime \prime})
\]

for \(g' \in G'\) and \(g'' \in G''\) (Proposition 2 of III, §9, No.~2, p.~542).

\begin{enumerate}
\def\labelenumi{\arabic{enumi})}
\setcounter{enumi}{5}
\tightlist
\item
  Let \((V,\pi)\) be a linear representation of G such that V is a finite-dimensional free K-module. We define a representation \(\rho\) of \(G\times G\) in \(\operatorname{End}_{K}(V)\) by the formula
\end{enumerate}

\[
\rho (g, g ^ {\prime}) (u) = \pi (g ^ {\prime}) \circ u \circ \pi (g ^ {- 1}).
\]

The K-module isomorphism \(\theta_{\mathrm{V}}: \mathrm{V}^{*} \otimes_{\mathrm{K}} \mathrm{V} \to \mathrm{End}_{\mathrm{K}}(\mathrm{V})\) of II, §4, No.~2, p.~271 is an isomorphism of representations from the external tensor product \(\pi^{\vee} \boxtimes \pi\) to \(\rho\) . The mapping \(g \mapsto \rho(1, g)\) (resp. \(g \mapsto \rho(g, 1)\) ) is a representation of G isomorphic to \(\pi^{\dim_{\mathrm{K}} \mathrm{V}}\) (resp. \((\pi^{\vee})^{\dim_{\mathrm{K}} \mathrm{V}}\) ).

Let L be a commutative K-algebra, and let \((\mathrm{M},\pi)\) be a linear representation of the group G. The group homomorphism \(\pi_{(\mathrm{L})}:\mathrm{G}\to\mathbf{GL}(\mathrm{M}_{(\mathrm{L})})\) defined by \(g\mapsto\mathrm{Id}_{\mathrm{L}}\otimes\pi(g)\) is a linear representation of G in the L-module \(\mathrm{M}_{(\mathrm{L})}\) , called the linear representation of G deduced from the representation \(\pi\) by extension of the ring of scalars K to L.

Suppose that K is a field and that L is a nonzero commutative K algebra. Let \((M,\pi)\) and \((M',\pi')\) be linear representations of G. The representations \(\pi\) and \(\pi'\) are isomorphic if and only if \(\pi_{(L)}\) and \(\pi'_{(L)}\) are (VIII, p.~37, Theorem 3).

Suppose, moreover, that the algebra L is an extension of K. Consider the rings \(\mathrm{R}_{\mathrm{K}}(\mathrm{G})\) and \(\mathrm{R}_{\mathrm{L}}(\mathrm{G})\) defined in Example 1 of VIII, p.~198. Extension of scalars defines a ring homomorphism

\[
u: \mathrm{R} _ {\mathrm{K}} (\mathrm{G}) \longrightarrow \mathrm{R} _ {\mathrm{L}} (\mathrm{G}).
\]

This homomorphism is injective, and an element \(\xi \in \mathrm{R}_{\mathrm{K}}(\mathrm{G})\) is effective if and only if \(u(\xi)\) is (VIII, p.~195, Theorem 1).

\subsection{Maschke's Theorem}

THEOREM 1. --- Suppose that the group G is finite. Let M be a K{[}G{]}-module, and let N be a K{[}G{]}-submodule of M. Suppose that N is a direct factor of the K-module M and that \textbar G\textbar{} is invertible in K. Then N is a direct factor of the K{[}G{]}-module M.

Let p be a K-linear projector in M with image N. We define an endomorphism q of the K-module M by setting

\[
q (x) = | \mathrm{G} | ^ {- 1} \sum_ {g \in \mathrm{G}} g p (g ^ {- 1} x)
\]

for every \(x \in M\) . Since N is stable under the action of G and p induces the identity on N, we see that q sends M into N and induces the identity on N.

The K-module M is therefore the direct sum of the image N of q and the kernel of q. We have \(g q(x) = q(gx)\) for every \(x \in M\) and \(g \in G\) , so that the kernel of q is a K{[}G{]}-submodule of M. This proves that N is a direct factor of the K{[}G{]}-module M.

COROLLARY 1 (Maschke). --- Suppose that the group G is finite and that K is a commutative field. The algebra K{[}G{]} is semisimple if and only if the element \textbar G\textbar{} of the field K is not zero.

Suppose \(|\mathrm{G}| \neq 0\) . By Theorem 1, every K{[}G{]}-submodule of \(\mathrm{K}[\mathrm{G}]_s\) is a direct factor. Hence K{[}G{]} is semisimple.

Conversely, suppose that \(|G|\) is zero, and denote by \(\varepsilon\) the element \(\sum_{g\in G} g\) of the center of K{[}G{]}. We have \(\varepsilon \neq 0\) but \(\varepsilon^{2} = |G|\varepsilon = 0\) , so K{[}G{]} is not semisimple (VIII, p.~153, Proposition 3, a) and VIII, p.~157, Remark 1).

COROLLARY 2. --- Suppose that the group G is finite and that \(|G|\) is invertible in K. An exact sequence of K{[}G{]}-modules splits if and only if it splits as an exact sequence of K-modules.

Given an exact sequence of K{[}G{]}-modules

\[
0 \longrightarrow \mathrm{M} ^ {\prime} \stackrel {f} {\longrightarrow} \mathrm{M} \longrightarrow \mathrm{M} ^ {\prime \prime} \longrightarrow 0,
\]

it suffices to apply Theorem 1 to the image of the morphism f.

COROLLARY 3. --- Suppose that the group G is finite and that \(|G|\) is invertible in K. A K{[}G{]}-module is projective if and only if it is projective as a K-module.

Let P be a K{[}G{]}-module. If P is a direct factor of a free K{[}G{]}-module M, then it is a fortiori a direct factor of the free K-module M. The converse follows from Corollary 2, in view of II, §2, No.~2, p.~231, Proposition 4, d).

COROLLARY 4. --- a) Suppose that the group G is finite and that K is a commutative field of characteristic 0. Two finite-dimensional linear representations of G are isomorphic if and only if they have the same character.

\begin{enumerate}
\def\labelenumi{\alph{enumi})}
\setcounter{enumi}{1}
\tightlist
\item
  Suppose that K is a perfect field of characteristic a prime number p and that the group G is finite, of cardinal prime to p.~The character of a finite-dimensional linear representation of G is zero if and only if this representation is isomorphic to the direct sum of p pairwise isomorphic representations.
\end{enumerate}

Under the assumptions of the corollary, every finite-dimensional linear representation of G is semisimple. The corollary then follows from the corollary of VIII, p.~384.

COROLLARY 5. --- Suppose that the group G is finite and that K is a commutative field in which \(|G| \neq 0\) . Let \(\pi\) and \(\pi'\) be finite-dimensional linear representations of G. Then \(\pi\) and \(\pi'\) are isomorphic if and only if for every g in G, the endomorphisms \(\pi(g)\) and \(\pi'(g)\) have the same characteristic polynomial.

This follows from Corollary 1 of VIII, p.~386.

\subsection{Induced and Coinduced Representations}

Let H be a subgroup of the group G.

If \((V,\pi)\) is a linear representation of G, then the restriction of \(\pi\) to H is a linear representation of H in V; we denote it by \(\operatorname{Res}_{\mathrm{H}}^{\mathrm{G}}(\pi)\) . The K{[}H{]}-module associated with \(\operatorname{Res}_{\mathrm{H}}^{\mathrm{G}}(\pi)\) is simply the module deduced from the K{[}G{]}-module V by restriction of scalars (II, §1, No.~13, p.~221). If V is a finite-dimensional free K-module, then the character of \(\operatorname{Res}_{\mathrm{H}}^{\mathrm{G}}(\pi)\) is the restriction of the character of \(\pi\) to H.

Let \((\mathrm{M},\sigma)\) be a representation of H.

We view K{[}G{]} as a (K{[}G{]}, K{[}H{]})-bimodule and M as a K{[}H{]}-module. The K{[}G{]}-module \(\mathcal{T}(M) = K[G] \otimes_{K[H]} M\) (VIII, p.~58) defines a linear representation of G, denoted by \(\operatorname{Ind}_{\mathrm{H}}^{\mathrm{G}}(\sigma)\) and called the representation of G induced by \(\sigma\) . If \((\mathrm{V}, \pi)\) is a linear representation of G, then the K{[}H{]}-module \(\mathcal{H}(\mathrm{V}) = \mathrm{Hom}_{\mathrm{K[G]}}(\mathrm{K[G]},\mathrm{V})\) can be identified with the K{[}H{]}-module corresponding to the representation \(\mathrm{Res}_{\mathrm{H}}^{\mathrm{G}}(\pi)\) . Consequently, the adjunction morphism (VIII, p.~59) gives a K-module isomorphism, called canonical, from \(\mathrm{Hom}_{\mathrm{H}}(\sigma ,\mathrm{Res}_{\mathrm{H}}^{\mathrm{G}}(\pi))\) to \(\mathrm{Hom}_{\mathrm{G}}(\mathrm{Ind}_{\mathrm{H}}^{\mathrm{G}}(\sigma),\pi)\) (``Frobenius reciprocity'').

We view K{[}G{]} as a (K{[}H{]}, K{[}G{]})-bimodule. The K{[}G{]}-module \(\mathcal{H}(M) = \operatorname{Hom}_{\mathrm{K}[\mathrm{H}]}(\mathrm{K}[\mathrm{G}], \mathrm{M})\) defines a representation of G, denoted by \(\operatorname{Coind}_{\mathrm{H}}^{\mathrm{G}}(\sigma)\) and called the representation of G coinduced by \(\sigma\) . If \((\mathrm{V}, \pi)\) is a linear representation of G, then the K{[}H{]}-module \(\mathcal{T}(\mathrm{V}) = \mathrm{K}[\mathrm{G}] \otimes_{\mathrm{K}[\mathrm{G}]} \mathrm{V}\) can be identified with the K{[}H{]}-module corresponding to the representation \(\operatorname{Res}_{\mathrm{H}}^{\mathrm{G}}(\pi)\) . Consequently, the adjunction morphism (loc. cit.) gives a K-module isomorphism, called canonical, from \(\operatorname{Hom}_{\mathrm{H}}(\operatorname{Res}_{\mathrm{H}}^{\mathrm{G}}(\pi), \sigma)\) to \(\operatorname{Hom}_{\mathrm{G}}(\pi, \operatorname{Coind}_{\mathrm{H}}^{\mathrm{G}}(\sigma))\) .

Let \(\varepsilon : \mathrm{K}[\mathrm{G}] \to \mathrm{K}[\mathrm{H}]\) be the K-module homomorphism characterized by the relations \(\varepsilon(h) = h\) if \(h \in \mathrm{H}\) and \(\varepsilon(g) = 0\) if \(g \in \mathrm{G} - \mathrm{H}\) . The mapping \(\varepsilon\) is a homomorphism of \((\mathrm{K}[\mathrm{H}], \mathrm{K}[\mathrm{H}])\) -bimodules. Let \((\mathrm{M}, \sigma)\) be a linear representation of H. The mapping \(v \mapsto v \circ \varepsilon\) from \(\mathrm{Hom}_{\mathrm{K}[\mathrm{H}]}(\mathrm{K}[\mathrm{H}], \mathrm{M})\) to \(\mathrm{Hom}_{\mathrm{K}[\mathrm{H}]}(\mathrm{K}[\mathrm{G}], \mathrm{M})\) is a homomorphism of \(\mathrm{K}[\mathrm{H}]\) -modules. By identifying M with \(\mathrm{Hom}_{\mathrm{K}[\mathrm{H}]}(\mathrm{K}[\mathrm{H}], \mathrm{M})\) , we obtain a homomorphism of \(\mathrm{K}[\mathrm{H}]\) -modules from M to \(\mathrm{Res}_{\mathrm{H}}^{\mathrm{G}}(\mathrm{Coind}_{\mathrm{H}}^{\mathrm{G}}(\sigma))\) . Frobenius reciprocity sends it to a homomorphism \(\iota\) of \(\mathrm{K}[\mathrm{G}]\) -modules from \(\mathrm{Ind}_{\mathrm{H}}^{\mathrm{G}}(\sigma)\) to \(\mathrm{Coind}_{\mathrm{H}}^{\mathrm{G}}(\sigma)\) characterized by the relation

\[
\iota (g \otimes m) (g ^ {\prime}) = \varepsilon (g ^ {\prime} g) m
\]

for \(g, g' \in G\) and \(m \in M\) . This homomorphism is called canonical.

PROPOSITION 2. --- Let H be a subgroup of G, and let \((\mathrm{M},\sigma)\) be a linear representation of H. The canonical homomorphism \(\iota:\operatorname{Ind}_{\mathrm{H}}^{\mathrm{G}}(\sigma)\to\operatorname{Coind}_{\mathrm{H}}^{\mathrm{G}}(\sigma)\) is injective. If the subgroup H has finite index in G, then \(\iota\) is an isomorphism of K{[}G{]}-modules.

Let \(S \subset G\) be a system of representatives of \(G/H\) . The family \((s)_{s \in S}\) is a basis of the right \(K[H]\) -module \(K[G]\) . It follows that the mapping \(M^{(S)} \to \operatorname{Ind}_{H}^{G}(\sigma)\) defined by \((m_s)_{s \in S} \mapsto \sum_{s \in S} s \otimes m_s\) is an isomorphism of \(K\) -modules. For all \(s, s' \in S\) and every \(m \in M\) , we have the relations

\[
\iota (s ^ {\prime} \otimes m) (s ^ {- 1}) = \left\{ \begin{array}{l} m \text {   if   } s = s ^ {\prime}, \\ 0 \text {   otherwise }. \end{array} \right.
\]

It follows that \(\iota'\) is injective.

Suppose that H has finite index in G. The set S is then finite; let \(\rho\) be the mapping from \(\operatorname{Coind}_{\mathrm{H}}^{\mathrm{G}}(\sigma)\) to \(\operatorname{Ind}_{\mathrm{H}}^{\mathrm{G}}(\sigma)\) given by \(u \mapsto \sum_{s \in S} s \otimes u(s^{-1})\) . It satisfies \((\iota \circ \rho(u))(s^{-1}) = u(s^{-1})\) for \(u \in \operatorname{Coind}_{\mathrm{H}}^{\mathrm{G}}(\sigma)\) and \(s \in S\) . Since the family \((s^{-1})_{s\in\mathrm{S}}\) is a basis of the left K{[}H{]}-module K{[}G{]}, the mapping \(\iota\circ\rho\) is the identity mapping. Consequently, \(\iota\) is bijective with inverse bijection \(\rho\) .

Let H be a subgroup of G of finite index. Let u be a central function on H; denote by \(u^{0}\) the function on G that extends u and that is zero at every point of G---H. Let S be a system of representatives of G/H. For \(g \in G\) , set

\[
\operatorname{Ind} _ {\mathrm{H}} ^ {\mathrm{G}} (u) (g) = \sum_ {s \in \mathrm{S}} u ^ {0} (s ^ {- 1} g s).\tag{4}
\]

Note that for all \(x, g \in G\) and every \(h \in H\) , we have \(u^{0}((xh)^{-1}gxh) = u^{0}(x^{-1}gx)\) . It follows that \(\operatorname{Ind}_{\mathrm{H}}^{\mathrm{G}}(u)\) is a central function on G that does not depend on the choice of S. We thus define a K-linear mapping \(Ind_{H}^{G}\) from \(\mathcal{Z}_{\mathrm{K}}(\mathrm{H})\) to \(\mathcal{Z}_{\mathrm{K}}(\mathrm{G})\) .

PROPOSITION 3. --- Let H be a subgroup of G of finite index. Let \((M,\sigma)\) be a linear representation of H. Suppose that M is a finite-dimensional free K-module. Denote by \((V,\pi)\) the representation of G induced by \(\sigma\) . Then the K-module V is free and finite-dimensional, and

\[
\chi_ {\pi} = \mathrm{Ind} _ {\mathrm{H}} ^ {\mathrm{G}} (\chi_ {\sigma}).\tag{5}
\]

Let S be a system of representatives of G/H. As we saw above, the linear mapping \(\mathrm{M}^{\mathrm{S}}\to\mathrm{Ind}_{\mathrm{H}}^{\mathrm{G}}(\mathrm{M})\) given by \((m_{s})_{s\in\mathrm{S}}\mapsto\sum_{s\in\mathrm{S}}s\otimes m_{s}\) is a K-module isomorphism. In particular, V is a finite-dimensional free K-module. For any \(g\in G\) , denote by \(M_{g}\) the image of M by the mapping \(m\mapsto g\otimes m\) . For all \(g,g'\in G\) , we have \(M_{g}=M_{g'}\) if and only if \(gH=g'H\) , and V is the direct sum of the submodules \(M_{s}\) for \(s\in S\) . For all \(g,g'\in G\) , we also have \(\pi(g)\mathrm{M}_{g'}=\mathrm{M}_{gg'}\) . For any \(g\in G\) , denote by \(S_{g}\) the set of elements s of S such that \(s^{-1}gs\in H\) . For all \(g,g'\in G\) such that \(g'^{-1}gg'\in H\) , the automorphism \(\pi(g)\) of V induces an automorphism \(\pi(g)_{g'}\) of \(M_{g'}\) , and we have

\[
\operatorname{Tr} (\pi (g)) = \sum_ {s \in \mathrm{S} _ {g}} \operatorname{Tr} (\pi (g) _ {s}) = \sum_ {s \in \mathrm{S} _ {g}} \operatorname{Tr} (\sigma (s ^ {- 1} g s)).
\]

The last assertion of the proposition follows.

\subsection{Representations and the Grothendieck Group}

In this subsection, we assume that K is a commutative field. Given a finite-dimensional linear representation \((M,\pi)\) of G, we denote by \([\pi]\) the class of the K{[}G{]}-module M in the Grothendieck ring \(\mathrm{R}_{\mathrm{K}}(\mathrm{G})\) (VIII, p.~198, Example 1).

By the definition of the laws of composition on the ring \(\mathrm{R}_{\mathrm{K}}(\mathrm{G})\) , if \(\pi\) and \(\pi'\) are finite-dimensional linear representations of G, then we have

\[
[ \pi \oplus \pi^ {\prime} ] = [ \pi ] + [ \pi^ {\prime} ], \quad [ \pi \otimes \pi^ {\prime} ] = [ \pi ] [ \pi^ {\prime} ].
\]

The unit element of the ring \(R_{K}(G)\) is the class of the unit representation of G (VIII, p.~399, Example 1).

The ring \(R_{K}(G)\) is a free Z-module with basis the set of classes of simple K{[}G{]}-modules of finite dimension over K. When \(|G|\) is invertible in K, two finite-dimensional representations are isomorphic if and only if they have the same class in \(R_{K}(G)\) (VIII, p.~190, Corollary and p.~401, Corollary 1).

For every finite-dimensional linear representation \((M,\pi)\) of G, the contragredient \(\pi^{\vee}\) is also finite-dimensional. The representations \(\pi\) and \((\pi^{\vee})^{\vee}\) are isomorphic. If \(\pi'\) is also a finite-dimensional linear representation of G, then the representations \((\pi\otimes\pi')^{\vee}\) and \(\pi^{\vee}\otimes\pi'^{\vee}\) are isomorphic. If

\[
0 \longrightarrow \mathrm{M} ^ {\prime} \longrightarrow \mathrm{M} \longrightarrow \mathrm{M} ^ {\prime \prime} \longrightarrow 0
\]

is an exact sequence of K{[}G{]}-modules of finite dimension over K, then the contragredient representations give an exact sequence of K{[}G{]}-modules

\[
0 \longrightarrow \mathrm{M} ^ {\prime \prime \vee} \longrightarrow \mathrm{M} ^ {\vee} \longrightarrow \mathrm{M} ^ {\prime \vee} \longrightarrow 0
\]

(II, §7, No.~5, p.~299, Corollary of Theorem 5). By the universal property of Grothendieck groups (VIII, p.~186, Proposition 4), there exists an automorphism \(c \mapsto c^{\vee}\) of the ring \(\mathrm{R}_{\mathrm{K}}(\mathrm{G})\) characterized by \([\pi]^{\vee} = [\pi^{\vee}]\) for every finite-dimensional linear representation \(\pi\) of \(\mathrm{G}\) ; we have \((c^{\vee})^{\vee} = c\) for every element \(c\) of \(\mathrm{R}_{\mathrm{K}}(\mathrm{G})\) .

Let H be a subgroup of G. Restriction preserves exact sequences. By the universal property of Grothendieck groups (loc. cit.), there consequently exists a group homomorphism, denoted by \(Res_{H}^{G}\) , from \(R_{K}(G)\) to \(R_{K}(H)\) characterized by the relation

\[
\mathrm{Res} _ {\mathrm{H}} ^ {\mathrm{G}} [ \pi ] = [ \mathrm{Res} _ {\mathrm{H}} ^ {\mathrm{G}} (\pi) ]\tag{6}
\]

for every finite-dimensional linear representation \(\pi\) of G; it is a ring homomorphism. If H has finite index in G, then by Proposition 14 of II, §7, No.~7, p.~306, we can define, analogously, a group homomorphism \(\mathrm{Ind}_{\mathrm{H}}^{\mathrm{G}}\) from \(\mathrm{R}_{\mathrm{K}}(\mathrm{H})\) to \(\mathrm{R}_{\mathrm{K}}(\mathrm{G})\) characterized by the relation

\[
\mathrm{Ind} _ {\mathrm{H}} ^ {\mathrm{G}} [ \sigma ] = [ \mathrm{Ind} _ {\mathrm{H}} ^ {\mathrm{G}} (\sigma) ]\tag{7}
\]

for every finite-dimensional linear representation \(\sigma\) of H.

By relations (1) of VIII, p.~399 and (3) of VIII, p.~400 and the universal property of Grothendieck groups mentioned above, there exists a ring homomorphism \(\Theta_{\mathrm{G}}\) from \(\mathrm{R}_{\mathrm{K}}(\mathrm{G})\) to the algebra \(\mathcal{Z}_{\mathrm{K}}(\mathrm{G})\) of central functions on \(\mathrm{G}\) , characterized by \(\Theta_{\mathrm{G}}([\pi]) = \chi_{\pi}\) for every finite-dimensional representation \(\pi\) of \(\mathrm{G}\) .

If H is a subgroup of G, then the homomorphisms \(\Theta_{G}\) and \(\Theta_{H}\) associated with the groups G and H are compatible with the operations \(Res_{H}^{G}\) and \(Ind_{H}^{G}\) (VIII, p.~402 and VIII, p.~404, Proposition 3).

Suppose that G is the product \(G^{\prime} \times G^{\prime\prime}\) of two groups. By VIII, p.~213, Remark 2, there exists a Z-linear mapping \(\kappa\) from \(\mathrm{R}_{\mathrm{K}}(\mathrm{G}^{\prime}) \otimes_{\mathbf{Z}} \mathrm{R}_{\mathrm{K}}(\mathrm{G}^{\prime\prime})\) to the group \(\mathrm{R}_{\mathrm{K}}(\mathrm{G}^{\prime} \times \mathrm{G}^{\prime\prime})\) characterized by the relation \(\kappa([\pi^{\prime}] \otimes [\pi^{\prime\prime}]) = [\pi^{\prime} \boxtimes \pi^{\prime\prime}]\) for \(\pi^{\prime}\) (resp. \(\pi^{\prime\prime}\) ) a finite-dimensional representation of \(G^{\prime}\) (resp. \(G^{\prime\prime}\) ) and \(\pi^{\prime} \boxtimes \pi^{\prime\prime}\) the external tensor product (VIII, p.~400, Example 5). It is a ring homomorphism. If the field K is algebraically closed, then the mapping \(\kappa\) is an isomorphism (VIII, p.~213, Remark 2).

Suppose that G is the product \(G' \times G''\) of two groups. Denote by \(\psi\) the homomorphism from \(\mathcal{Z}_{\mathrm{K}}(\mathrm{G}') \otimes_{\mathrm{K}} \mathcal{Z}_{\mathrm{K}}(\mathrm{G}'' )\) to \(\mathcal{Z}_{\mathrm{K}}(\mathrm{G})\) that sends \(f' \otimes f''\) to the function \((g', g'') \mapsto f'(g')f''(g'')\) . The following diagram commutes:

\[
\mathrm{R} _ {\mathrm{K}} (\mathrm{G} ^ {\prime}) \otimes_ {\mathbf {Z}} \mathrm{R} _ {\mathrm{K}} (\mathrm{G} ^ {\prime \prime}) \xrightarrow {\kappa} \mathrm{R} _ {\mathrm{K}} (\mathrm{G})
\]

\[
\bigg \downarrow \Theta_ {\mathrm{G} ^ {\prime}} \otimes \Theta_ {\mathrm{G} ^ {\prime \prime}} \qquad \qquad \bigg \downarrow \Theta_ {\mathrm{G}}
\]

\[
\mathcal {Z} _ {\mathrm{K}} (\mathrm{G} ^ {\prime}) \otimes_ {\mathrm{K}} \mathcal {Z} _ {\mathrm{K}} (\mathrm{G} ^ {\prime \prime}) \xrightarrow {\psi} \mathcal {Z} _ {\mathrm{K}} (\mathrm{G}).
\]

\subsection{Fourier Inversion Formula}

For the remainder of this section, we assume that the group G is finite and that K is an algebraically closed field whose characteristic does not divide the order of G, so that the element \(|G|\) of K is not zero.

The algebra K{[}G{]} is semisimple (Maschke's theorem) and finite-dimensional. Denote by \(\widehat{G}\) the set of classes of simple K{[}G{]}-modules. For every \(\lambda\in\widehat{G}\) , choose a linear representation \((\mathrm{V}_{\lambda},\pi_{\lambda})\) of G such that the associated K{[}G{]}-module has class \(\lambda\) . The set \(\widehat{G}\) is finite, and the vector spaces \(V_{\lambda}\) are finite-dimensional (VIII, p.~141, Example). For any \(\lambda\in\widehat{G}\) , denote by \(d_{\lambda}\) the degree of the representation \(\pi_{\lambda}\) , that is, the dimension of the K-vector space \(V_{\lambda}\) , and denote its character by \(\chi_{\lambda}\) .

Denote by \(\mathrm{F}(\widehat{\mathrm{G}})\) the product algebra \(\prod_{\lambda\in\widehat{G}}\operatorname{End}_{\mathrm{K}}(\mathrm{V}_{\lambda})\) and by \(\overline{F}\) the mapping from K{[}G{]} to \(\mathrm{F}(\widehat{\mathrm{G}})\) defined by \(\overline{\mathcal{F}}(a)=(\pi_{\lambda}(a))_{\lambda\in\widehat{G}}\) . Since the field K is algebraically closed, the mapping \(\overline{F}\) is an algebra isomorphism (loc. cit.).

For every \(\lambda \in \widehat{\mathbf{G}}\) , the dimension of the algebra \(\mathrm{End}_{\mathrm{K}}(\mathrm{V}_{\lambda})\) is \(d_{\lambda}^{2}\) ; that of the algebra \(\mathrm{K}[\mathrm{G}]\) is \(\mathrm{Card}(\mathrm{G})\) . We therefore have the relation

\[
\mathrm{Card(G)} = \sum_ {\lambda \in \widehat {\mathrm{G}}} d _ {\lambda} ^ {2}.\tag{8}
\]

Denote by \(\tau\) the trace in the algebra K{[}G{]}; by definition, the trace \(\tau(a)\) of an element \(a\) of K{[}G{]} is the trace of the endomorphism \(x \mapsto ax\) of K{[}G{]} (III, §7, No.~7, p.~515). Let \(a = \sum_{g \in G} a_g g\) be an element of K{[}G{]}; by formula (2), we have \(\tau(ag^{-1}) = |G| a_g\) for every \(g \in G\) , and therefore the relation

\[
a = | \mathrm{G} | ^ {- 1} \sum_ {g \in \mathrm{G}} \tau (a g ^ {- 1}) g.\tag{9}
\]

Denote by \(\widehat{\tau}\) the trace in the algebra \(F(\widehat{G})\) . Let \(A = (A_{\lambda})_{\lambda \in \widehat{G}}\) be an element of \(F(\widehat{G})\) ; we have (cf.~III, §9, No.~3, p.~545, Example 3)

\[
\widehat {\tau} (\mathrm{A}) = \sum_ {\lambda \in \widehat {\mathrm{G}}} d _ {\lambda} \operatorname{Tr} (\mathrm{A} _ {\lambda}).\tag{10}
\]

Since the mapping \(\overline{F}\) is a K-algebra isomorphism, we have \(\widehat{\tau}\circ\overline{F}=\tau\) , and therefore

\[
\tau (a) = \widehat {\tau} (\overline {{\mathcal {F}}} (a)) = \widehat {\tau} \big ((\pi_ {\lambda} (a)) _ {\lambda \in \widehat {\mathrm{G}}} \big) = \sum_ {\lambda \in \widehat {\mathrm{G}}} d _ {\lambda} \operatorname{Tr} (\pi_ {\lambda} (a)) = \sum_ {\lambda \in \widehat {\mathrm{G}}} d _ {\lambda} \operatorname{Tr} _ {\lambda} (a)
\]

for every \(a\in \mathrm{K}[\mathrm{G}]\) , that is

\[
\tau = \sum_ {\lambda \in \widehat {\mathrm{G}}} d _ {\lambda} \operatorname{Tr} _ {\lambda}.\tag{11}
\]

Therefore, by (2) of VIII, p.~399, for \(g \in \mathbf{G}\) , we have

\[
\sum_ {\lambda \in \widehat {G}} d _ {\lambda}   \chi_ {\lambda} (g) = \left\{ \begin{array}{l l} | G | & \text { if   } g \text {   is   the   identity   element }, \\ 0 & \text { otherwise }. \end{array} \right.\tag{12}
\]

For \(a \in \mathrm{K}[\mathrm{G}]\) , relation (9) takes on the following form:

\[
a = | \mathrm{G} | ^ {- 1} \sum_ {g \in \mathrm{G}} \sum_ {\lambda \in \widehat {\mathrm{G}}} d _ {\lambda} \left. \operatorname{Tr} (\pi_ {\lambda} (a) \pi_ {\lambda} (g ^ {- 1})) g; \right.\tag{13}
\]

it follows that for every element \(A = (A_{\lambda})_{\lambda \in \widehat{G}}\) of \(F(\widehat{G})\) , we have

\[
\overline {{\mathcal {F}}} ^ {- 1} (\mathrm{A}) = | \mathrm{G} | ^ {- 1} \sum_ {g \in \mathrm{G}} \sum_ {\lambda \in \widehat {\mathrm{G}}} d _ {\lambda} \operatorname{Tr} (\mathrm{A} _ {\lambda} \pi_ {\lambda} (g ^ {- 1})) g\tag{14}
\]

(``Fourier inversion formula'').

For \(\mu \in \widehat{\mathrm{G}}\) , denote by \(j_{\mu}:\mathrm{End}_{\mathrm{K}}(\mathrm{V}_{\mu})\longrightarrow \prod_{\lambda \in \widehat{\mathrm{G}}} \mathrm{End}_{\mathrm{K}}(\mathrm{V}_{\lambda})\) the mapping such that \(j_{\mu}(u) = (v_{\lambda})\) , where \(v_{\lambda} = 0\) if \(\lambda \neq \mu\) and \(v_{\mu} = u\) . By formula (14), we have

\[
\overline {{\mathcal {F}}} ^ {- 1} (j _ {\mu} (u)) = | \mathrm{G} | ^ {- 1} d _ {\mu} \sum_ {g \in \mathrm{G}} \mathrm{Tr} (u \pi_ {\mu} (g ^ {- 1})) g.\tag{15}
\]

The center of the algebra \(\mathrm{F}(\widehat{\mathrm{G}})=\prod_{\lambda\in\widehat{\mathrm{G}}} \mathrm{End}_{\mathrm{K}}(\mathrm{V}_{\lambda})\) consists of the families \((a_{\lambda}1_{\mathrm{V}_{\lambda}})_{\lambda\in\widehat{\mathrm{G}}}\) , where \((a_{\lambda})\) is a family of elements of K. It is the image by \(\overline{{\mathcal{F}}}\) of the center of the algebra K{[}G{]}. That center therefore has a basis \((e_{\lambda})_{\lambda\in\widehat{\mathrm{G}}}\) characterized by the relation

\[
\pi_ {\lambda} (e _ {\mu}) = \delta_ {\lambda \mu} 1 _ {V _ {\lambda}}\tag{16}
\]

for \(\lambda, \mu \in \widehat{\mathbf{G}}\) , where \(\delta_{\lambda \mu}\) is the Kronecker delta function. By formula (15), we have, for every \(\mu \in \widehat{\mathbf{G}}\) ,

\[
e _ {\mu} = | \mathrm{G} | ^ {- 1} d _ {\mu} \sum_ {g \in \mathrm{G}} \chi_ {\mu} (g ^ {- 1}) g.\tag{17}
\]

These elements satisfy the relations

\[
\sum_ {\lambda \in \widehat {\mathsf {G}}} e _ {\lambda} = 1, \quad e _ {\mu} ^ {2} = e _ {\mu}, \quad \mathrm{and} \quad e _ {\mu} e _ {\nu} = 0\tag{18}
\]

for all \(\mu, \nu \in \widehat{\mathrm{G}}\) such that \(\mu \neq \nu\) ; they are the indecomposable idempotents of \(\mathrm{Z}(\mathrm{K}[\mathrm{G}])\) (VIII, p.~147, Proposition 15)

Remark. --- Let \((V, \pi)\) be a linear representation of G. By Maschke's theorem, the K{[}G{]}-module V is semisimple. For any \(\lambda \in \widehat{G}\) , denote by \(V^{\lambda}\) the isotypical component of type \(\lambda\) of the K{[}G{]}-module V; we have \(V = \bigoplus_{\lambda \in \widehat{G}} V^{\lambda}\) . By Proposition 15 of VIII, p.~147 and formula (17), the projector of V with image \(V^{\lambda}\) associated with this decomposition of V is equal to

\[
\pi (e _ {\lambda}) = | \mathrm{G} | ^ {- 1} d _ {\lambda} \sum_ {g \in \mathrm{G}} \chi_ {\lambda} (g ^ {- 1}) \pi (g).\tag{19}
\]

Applying this formula to \(\pi = \pi_{\lambda}\) , we obtain that the element \(d_{\lambda} \cdot 1\) of K is not zero. We will see further on (VIII, p.~420, Corollary 2) that \(d_{\mu}\) divides the cardinal of G.

Let \(\lambda\) be an element of \(\widehat{\mathbf{G}}\) . The mapping \(\pi_{\lambda}\) from K{[}G{]} to \(\mathrm{End}_{\mathrm{K}}(\mathrm{V}_{\lambda})\) induces an isomorphism of (K{[}G{]}, K{[}G{]})-bimodules from \(e_{\lambda}\mathrm{K}[\mathrm{G}]\) to \(\mathrm{End}_{\mathrm{K}}(\mathrm{V}_{\lambda})\) . By Proposition 15 of VIII, p.~147, the K{[}G{]}-submodule \(e_{\lambda}\mathrm{K}[\mathrm{G}]\) is the isotypical component of K{[}G{]} of type \(\lambda\) . It is also the isotypical component of type \(\lambda^{\vee}\) for the right regular representation of G (VIII, p.~143, Proposition 11) as well as the isotypical component of type \(\lambda \times \lambda^{\vee}\) for the biregular representation (VIII, p.~381, Proposition 4).

\subsection{Schur Orthogonality Relations}

We keep the notation of the previous subsection. Let \(\lambda\) be an element of \(\widehat{G}\) , and let \(u\) and \(v\) be elements of \(\mathrm{End}_{\mathrm{K}}(\mathrm{V}_{\lambda})\) ; by formula (10), we have

\[
\hat {\tau} (j _ {\lambda} (u) j _ {\lambda} (v)) = d _ {\lambda} \operatorname{Tr} (u v).
\]

Since \(\tau = \hat{\tau} \circ \overline{\mathcal{F}}\) , we deduce from formulas (2) and (15) the relation

\[
d _ {\lambda} ^ {2} | \mathrm{G} | ^ {- 1} \sum_ {g \in \mathrm{G}} \operatorname{Tr} (u \pi_ {\lambda} (g)) \operatorname{Tr} (v \pi_ {\lambda} (g ^ {- 1})) = d _ {\lambda} \operatorname{Tr} (u v).
\]

We observed previously that \(d_{\lambda}.1 \neq 0\) in the field K; it follows that

\[
| \mathrm{G} | ^ {- 1} \sum_ {g \in \mathrm{G}} \operatorname{Tr} (u \pi_ {\lambda} (g)) \operatorname{Tr} (v \pi_ {\lambda} (g ^ {- 1})) = d _ {\lambda} ^ {- 1} \operatorname{Tr} (u v).\tag{20}
\]

Let us apply this relation to the case when u and v have rank \(\leqslant 1\) ; we obtain

\[
| \mathrm{G} | ^ {- 1} \sum_ {g \in \mathrm{G}} \left\langle x ^ {*}, \pi_ {\lambda} (g) x \right\rangle \left\langle y ^ {*}, \pi_ {\lambda} \left(g ^ {- 1}\right) y \right\rangle = d _ {\lambda} ^ {- 1} \left\langle x ^ {*}, y \right\rangle \left\langle y ^ {*}, x \right\rangle\tag{21}
\]

for \(x, y\) in \(\mathrm{V}_{\lambda}\) and \(x^{*}, y^{*}\) in the dual \(\mathrm{V}_{\lambda}^{*}\) of \(\mathrm{V}_{\lambda}\) .

For every \(\lambda \in \widehat{\mathbf{G}}\) , let \((e_{\lambda,j})_{1\leqslant j\leqslant d_{\lambda}}\) be a basis of \(\mathrm{V}_{\lambda}\) ; denote by \((\pi_{ij}^{\lambda}(g))\) the matrix of the endomorphism \(\pi_{\lambda}(g)\) of \(\mathrm{V}_{\lambda}\) with respect to this basis. If we denote by \((e_{\lambda,i}^{*})_{1\leqslant i\leqslant d_{\lambda}}\) the basis of \(\mathrm{V}_{\lambda}^{*}\) dual to \((e_{\lambda,j})\) , then we have \(\pi_{ij}^{\lambda}(g) = \langle e_{\lambda,i}^{*}, \pi_{\lambda}(g)e_{\lambda,j}\rangle\) , and therefore

\[
| \mathrm{G} | ^ {- 1} \sum_ {g \in \mathrm{G}} \pi_ {i j} ^ {\lambda} (g) \pi_ {k \ell} ^ {\lambda} (g ^ {- 1}) = d _ {\lambda} ^ {- 1} \delta_ {i \ell} \delta_ {j k}.\tag{22}
\]

Now let \(\lambda\) and \(\mu\) be two distinct elements of \(\widehat{\mathbf{G}}\) , and let \(u \in \mathrm{End}_{\mathrm{K}}(\mathrm{V}_{\lambda})\) and \(v \in \mathrm{End}_{\mathrm{K}}(\mathrm{V}_{\mu})\) . Again by relation (15), we have

\[
\sum_ {g \in \mathrm{G}} \operatorname{Tr} (u \pi_ {\lambda} (g)) \operatorname{Tr} (v \pi_ {\mu} (g ^ {- 1})) = 0.\tag{23}
\]

As above, we deduce that

\[
\sum_ {g \in \mathrm{G}} \langle x ^ {*}, \pi_ {\lambda} (g) x \rangle \langle y ^ {*}, \pi_ {\mu} (g ^ {- 1}) y \rangle = 0\tag{24}
\]

for \(x\in \mathrm{V}_{\lambda},x^{*}\in \mathrm{V}_{\lambda}^{*},y\in \mathrm{V}_{\mu}\) , and \(y^{*}\in \mathrm{V}_{\mu}^{*}\) . We also have

\[
\sum_ {g \in \mathrm{G}} \pi_ {i j} ^ {\lambda} (g) \pi_ {k \ell} ^ {\mu} (g ^ {- 1}) = 0\tag{25}
\]

for \(i,j\) in \([1,d_{\lambda}]\) and \(k,\ell\) in \([1,d_{\mu}]\) .

Relations (20) through (25) are known as the Schur orthogonality relations.

Remark. --- We identify the algebra \(\operatorname{End}_{\mathrm{K}}(\mathrm{V}_{\lambda})\) with the matrix algebra \(\mathbf{M}_{d_{\lambda}}(\mathrm{K})\) via the basis \((e_{\lambda,j})\) of \(V_{\lambda}\) . The mapping \(\overline{F}^{-1}\) is an isomorphism from the algebra \(\prod_{\lambda}\mathbf{M}_{d_{\lambda}}(\mathrm{K})\) to the algebra K{[}G{]}. For \(\mu\in\widehat{G}\) , denote by \(E_{ij}^{\mu}\) the element of \(\prod_{\lambda}\mathbf{M}_{d_{\lambda}}(\mathrm{K})\) whose component of index \(\mu\) is the matrix unit \(E_{ij}\) of \(\mathbf{M}_{d_{\mu}}(\mathrm{K})\) (II, §10, No.~3, p.~341) and whose other components are zero; set \(u_{ij}^{\mu}=\overline{\mathcal{F}}^{-1}(\mathrm{E}_{ij}^{\mu})\) . The family of elements \(u_{ij}^{\lambda}\) , for \(\lambda\in\widehat{G}\) , \(1\leqslant i\leqslant d_{\lambda}\) , \(1\leqslant j\leqslant d_{\lambda}\) , is a basis of the algebra K{[}G{]}; the multiplication table is

\[
u _ {i j} ^ {\lambda} u _ {k \ell} ^ {\mu} = \delta_ {\lambda \mu} \delta_ {j k} u _ {i \ell} ^ {\lambda}.\tag{26}
\]

Moreover, by formula (15), we have

\[
u _ {i j} ^ {\lambda} = | \mathrm{G} | ^ {- 1} d _ {\lambda} \sum_ {g \in \mathrm{G}} \pi_ {j i} ^ {\lambda} (g ^ {- 1}) g.\tag{27}
\]

\subsection{Orthogonality Relation for Characters}

We keep the notation of Subsections 5 and 6. Recall that \(Z(K[G])\) consists of the central functions.

We define a symmetric bilinear mapping from \(K[G] \times K[G]\) to K by the formula

\[
\langle f, f ^ {\prime} \rangle_ {\mathrm{G}} = | \mathrm{G} | ^ {- 1} \sum_ {g \in \mathrm{G}} f _ {g} f _ {g ^ {- 1}} ^ {\prime}\tag{28}
\]

for every \(f = \sum f_{g}g\) and \(f' = \sum f_{g}'g\) belonging to K{[}G{]}. We have \(\langle f, f' \rangle_{\mathrm{G}} = |\mathrm{G}|^{-2}\tau (ff')\) .

PROPOSITION 4 (Orthogonality relation for characters)

For \(\lambda\) and \(\mu\) in \(\hat{\mathrm{G}}\) , we have \(\langle \chi_{\lambda},\chi_{\mu}\rangle_{\mathrm{G}} = \delta_{\lambda \mu}\) .

This is the specific case of relations (20) and (23) when the endomorphisms u and v are taken to be the identity.

COROLLARY. --- Let \(\pi\) and \(\pi'\) be finite-dimensional linear representations of G. In the field K, we have

\[
\langle \chi_ {\pi}, \chi_ {\pi^ {\prime}} \rangle_ {\mathrm{G}} = (\dim_ {\mathrm{K}} \operatorname{Hom} _ {\mathrm{G}} (\pi , \pi^ {\prime})) \cdot 1.\tag{29}
\]

We first suppose that \(\pi\) and \(\pi'\) are simple representations. The vector space \(\mathrm{Hom}_{\mathrm{G}}(\pi, \pi')\) has dimension 1 or 0 according to whether or not \(\pi\) and \(\pi'\) are isomorphic (Schur's lemma, VIII, p.~47, Proposition 2). In this case, formula (29) follows from Proposition 4.

In the general case, the representation \(\pi\) (resp. \(\pi'\) ) is the direct sum of simple representations \(\pi_{1},\ldots,\pi_{m}\) (resp. \(\pi_{1}',\ldots,\pi_{n}'\) ). The space \(\operatorname{Hom}_{\mathrm{G}}(\pi,\pi')\) is isomorphic to the direct sum of the spaces \(\operatorname{Hom}_{\mathrm{G}}(\pi_{i},\pi_{j}')\) for \(1\leqslant i\leqslant m\) , \(1\leqslant j\leqslant n\) , and we have

\[
\chi_ {\pi} = \chi_ {\pi_ {1}} + \dots + \chi_ {\pi_ {m}}, \quad \chi_ {\pi^ {\prime}} = \chi_ {\pi_ {1} ^ {\prime}} + \dots + \chi_ {\pi_ {n} ^ {\prime}}.
\]

By linearity, the proof of formula \((29)\) is reduced to the case of simple representations.

\subsection{Central Functions on a Finite Group}

PROPOSITION 5. --- The family \((\chi_{\lambda})_{\lambda\in\widehat{\mathbf{G}}}\) is a basis of the vector space of central functions. The number of classes of simple linear representations of G is equal to the number of conjugacy classes of G.

For \(a = \sum_{g \in G} a_g g \in K[G]\) , write \(a^\vee = \sum_{g \in G} a_{g^{-1}} g\) ; the mapping \(a \mapsto a^\vee\) is an involutive antiautomorphism of the algebra K{[}G{]}. By formula (17), we have \(e_\lambda = |G|^{-1} d_\lambda \chi_\lambda^\vee\) for every \(\lambda \in \widehat{G}\) . The proposition then follows from the fact that the family \((e_\lambda)_{\lambda \in \widehat{G}}\) is a basis of the center of K{[}G{]}.

Denote by C the set of conjugacy classes of G. Let C be an element of C; the image \(C^{-1}\) of C by the mapping \(g \mapsto g^{-1}\) is a conjugacy class. The commutants of the elements of C are mutually conjugate subgroups of G; their cardinal \(d(\mathrm{C})\) satisfies

\[
\operatorname{Card} (\mathrm{G}) = \operatorname{Card} (\mathrm{C}) d (\mathrm{C}).\tag{30}
\]

In particular, we have \(d(\mathrm{C}) \cdot 1 \neq 0\) in the field K.

Let f be a central function on G. For any conjugacy class C, denote by \(f(\mathrm{C})\) the common value of the \(f(x)\) for \(x \in C\) . With this notation, the orthogonality relation for characters (VIII, p.~410, Proposition 4) can be written as

\[
\sum_ {\mathrm{C} \in \mathcal {C}} \chi_ {\lambda} (\mathrm{C} ^ {- 1}) \chi_ {\mu} (\mathrm{C}) d (\mathrm{C}) ^ {- 1} = \delta_ {\lambda \mu}\tag{31}
\]

for \(\lambda\) and \(\mu\) in \(\widehat{G}\) .

Denote by \(A\) the matrix of type \(\widehat{\mathbf{G}} \times \mathcal{C}\) with entries \(\chi_{\lambda}(\mathrm{C})\) and by \(B\) the matrix of type \(\mathcal{C} \times \widehat{\mathbf{G}}\) with entries \(\chi_{\lambda}(\mathrm{C}^{-1}) d(\mathrm{C})^{-1}\) . The sets \(\widehat{\mathbf{G}}\) and \(\mathcal{C}\) have the same cardinal (Proposition 5); relation (31) expresses the fact that the matrix product \(AB\) is the matrix unit of type \(\widehat{\mathbf{G}} \times \widehat{\mathbf{G}}\) . By Proposition 11 of II, §10, No.~12, p.~360, the matrix product \(BA\) is the matrix unit of type \(\mathcal{C} \times \mathcal{C}\) ; in other words, we have the relation

\[
\sum_ {\lambda \in \widehat {\mathrm{G}}} \chi_ {\lambda} (\mathrm{C} ^ {- 1}) \chi_ {\lambda} (\mathrm{C} ^ {\prime}) = d (\mathrm{C}) \delta_ {\mathrm{CC} ^ {\prime}}\tag{32}
\]

for C and C' in C (sometimes called the ``second orthogonality relation for characters'').

Let H be a subgroup of G. Let us note that the integer \(\text{Card(H)}\) divides \(\text{Card(G)}\) and that \(|G|\) is not zero in K; we therefore have \(|H| \neq 0\) in K.

Denote by \(Res_{H}^{G}\) the linear mapping from \(Z(K[G])\) to \(Z(K[H])\) that sends a central function on G to its restriction to H. We have seen that if \(\chi_{\pi}\) is the character of a finite-dimensional representation \(\pi\) of G, then \(\mathrm{Res}_{\mathrm{H}}^{\mathrm{G}}(\chi_{\pi})\) is the character of the representation \(\mathrm{Res}_{\mathrm{H}}^{\mathrm{G}}(\pi)\) of H.

PROPOSITION 6. --- Let f be a central function on G and u be a central function on H. We have

\[
\langle \mathrm{Ind} _ {\mathrm{H}} ^ {\mathrm{G}} (u), f \rangle_ {\mathrm{G}} = \langle u, \mathrm{Res} _ {\mathrm{H}} ^ {\mathrm{G}} (f) \rangle_ {\mathrm{H}}.\tag{33}
\]

The characters of the simple representations of G form a basis of \(Z(\mathrm{K}[\mathrm{G}])\) (VIII, p.~411, Proposition 5), and the same holds for H. It therefore suffices to establish (33) in the case when f is the character \(\chi_{\pi}\) of a simple representation \(\pi\) of G and u is the character \(\chi_{\sigma}\) of a simple representation \(\sigma\) of H. In this case, \(\operatorname{Ind}_{\mathrm{H}}^{\mathrm{G}}(u)\) is the character of the representation \(\operatorname{Ind}_{\mathrm{H}}^{\mathrm{G}}(\sigma)\) of G, and, by VIII, p.~410, Corollary, we have

\[
\langle \mathrm{Ind} _ {\mathrm{H}} ^ {\mathrm{G}} (u), f \rangle_ {\mathrm{G}} = (\dim_ {\mathrm{K}} \mathrm{Hom} _ {\mathrm{G}} (\mathrm{Ind} _ {\mathrm{H}} ^ {\mathrm{G}} (\sigma), \pi)) \cdot 1.
\]

We show the relation

\[
\langle u, \mathrm{Res} _ {\mathrm{H}} ^ {\mathrm{G}} (f) \rangle_ {\mathrm{H}} = (\dim_ {\mathrm{K}} \mathrm{Hom} _ {\mathrm{H}} (\sigma , \mathrm{Res} _ {\mathrm{H}} ^ {\mathrm{G}} (\pi))) \cdot 1
\]

likewise, and equality (33) follows by Frobenius reciprocity.

PROPOSITION 7. --- Let \(f\) be a mapping from \(G\) to \(K\) . The following assertions are equivalent:

\begin{enumerate}
\def\labelenumi{(\roman{enumi})}
\item
  There are an element \(\lambda\) of \(\widehat{\mathbf{G}}\) and an element \(a\) of \(\mathrm{K}\) such that \(f = a\chi_{\lambda}\) .
\item
  For every pair \((g,g^{\prime})\) of elements of \(\mathbf{G}\) , we have
\end{enumerate}

\[
f (g) f (g ^ {\prime}) = | \mathrm{G} | ^ {- 1} f (1) \sum_ {h \in \mathrm{G}} f (h g h ^ {- 1} g ^ {\prime}).\tag{34}
\]

Let \(\lambda\) be an element of \(\widehat{\mathbf{G}}\) ; for every endomorphism \(u\) of \(\mathrm{V}_{\lambda}\) , set

\[
u ^ {\natural} = | \mathrm{G} | ^ {- 1} \sum_ {h \in \mathrm{G}} \pi_ {\lambda} (h) u \pi_ {\lambda} (h ^ {- 1}).\tag{35}
\]

The endomorphism \(u^{\natural}\) of \(V_{\lambda}\) is K{[}G{]}-linear. By Schur's lemma (VIII, p.~47, Theorem 1), \(u^{\natural}\) is a homothety. Since \(u\) and \(u^{\natural}\) have the same trace, we therefore have

\[
u ^ {\natural} = d _ {\lambda} ^ {- 1} \mathrm{Tr} (u) 1 _ {\mathrm{V} _ {\lambda}}.
\]

Let \(u\) and \(v\) be endomorphisms of \(V_{\lambda}\) ; it follows that

\[
\operatorname{Tr} (u) \operatorname{Tr} (v) = d _ {\lambda} \operatorname{Tr} (u ^ {\natural} v) = d _ {\lambda} | G | ^ {- 1} \sum_ {h \in G} \operatorname{Tr} (\pi_ {\lambda} (h) u \pi_ {\lambda} (h ^ {- 1}) v).\tag{36}
\]

Take \(u = \pi_{\lambda}(g)\) and \(v = \pi_{\lambda}(g')\) in formula (36); relation (34) for \(f = \chi_{\lambda}\) follows.

Conversely, suppose that assertion (ii) holds. If we have \(f(1)=0\) , then it follows that \(f(g)f(g')=0\) for every pair \((g,g')\) of elements of G, and therefore f=0. We can therefore assume \(f(1)\neq0\) . Taking \(g'=1\) in (34), we obtain the relation

\[
f (g) = | \mathrm{G} | ^ {- 1} \sum_ {h \in \mathrm{G}} f (h g h ^ {- 1})
\]

for every \(g \in G\) , which implies that f is a central function. By Proposition 5, there exists a family \((a_{\lambda})\) of elements of K such that \(f = \sum_{\lambda \in \widehat{G}} a_{\lambda} \chi_{\lambda}\) . Let us replace f with this expression in formula (34); taking into account that each of the characters \(\chi_{\lambda}\) also satisfies this relation, we find

\[
\sum_ {\lambda , \mu \in \widehat {G}} a _ {\lambda}   a _ {\mu}   \chi_ {\lambda} (g)   \chi_ {\mu} (g ^ {\prime}) = \sum_ {\lambda \in \widehat {G}} a _ {\lambda}   d _ {\lambda} ^ {- 1}   f (1)   \chi_ {\lambda} (g)   \chi_ {\lambda} (g ^ {\prime})\tag{37}
\]

for \(g, g' \in \mathbf{G}\) . This relation can also be written as

\[
\sum_ {\lambda , \mu} (a _ {\lambda} a _ {\mu} - \delta_ {\lambda \mu} a _ {\lambda} d _ {\lambda} ^ {- 1} f (1)) \chi_ {\lambda} (g) \chi_ {\mu} (g ^ {\prime}) = 0\tag{38}
\]

for \(g, g' \in \mathrm{G}\) . Now, the functions \(\chi_{\lambda}\) , for \(\lambda \in \widehat{\mathrm{G}}\) , are linearly independent (Proposition 5 of VIII, p.~411); it follows that

\[
a _ {\lambda} a _ {\mu} = \delta_ {\lambda \mu} a _ {\lambda} d _ {\lambda} ^ {- 1} f (1)
\]

for \(\lambda, \mu \in \widehat{\mathbf{G}}\) . In particular, \(a_{\lambda}a_{\mu} = 0\) whenever \(\lambda \neq \mu\) . Consequently, there exists at most one element \(\lambda\) of \(\widehat{\mathbf{G}}\) such that \(a_{\lambda} \neq 0\) , and we have \(f = a_{\lambda}\chi_{\lambda}\) , and therefore (i).

\subsection{The Case of Abelian Groups}

In this subsection, we assume that the group \(\mathrm{G}\) is abelian.

By Schur's lemma (VIII, p.~48, Corollary 1), every simple representation of G has dimension 1. Let \((\mathrm{M},\pi)\) be such a representation and \(\chi\) be its character; for every \(g\in G\) and \(x\in M\) , we have \(\pi(g)(x)=\chi(g)x\) . Consequently, the character \(\chi\) is a homomorphism from G to the multiplicative group \(K^{*}\) of K. Conversely, every homomorphism from G to \(K^{*}\) is the character of a representation of G of degree 1. So the set \(\widehat{G}\) of classes of simple K{[}G{]}-modules can be identified with the set \(\operatorname{Hom}(G,K^{*})\) of homomorphisms from G to \(K^{*}\) . We deduce from this an abelian group structure on \(\widehat{G}\) ; the product in \(\widehat{G}\) corresponds to the tensor product of representations. The groups G and \(\widehat{G}\) have the same cardinal by Proposition 5 of VIII, p.~411. Every function on G is central, and \(\widehat{G}\) is a basis of the vector space of mappings from G to K (loc. cit.). Because of the orthogonality relation for characters, such a mapping f has the following unique decomposition with respect to the basis \(\widehat{G}\) :

\[
f = \sum_ {\chi \in \widehat {\mathrm{G}}} \langle \chi , f \rangle_ {\mathrm{G}} \chi .\tag{39}
\]

For mappings \(f\) and \(f'\) from \(G\) to \(K\) , we have the relation

\[
\langle f, f ^ {\prime} \rangle_ {\mathrm{G}} = \sum_ {\chi \in \widehat {\mathrm{G}}} \langle \chi , f \rangle_ {\mathrm{G}} \langle \chi , f ^ {\prime} \rangle_ {\mathrm{G}}.\tag{40}
\]

Let \((\mathrm{V},\pi)\) be a linear representation of G. For any \(\chi\in\widehat{\mathrm{G}}\) , denote by \(V^{\chi}\) the subspace of V consisting of the vectors v such that \(\pi(g)(v)=\chi(g)v\) for every \(g\in\mathrm{G}\) . The space \(V^{\chi}\) is the isotypical component of type \(\chi\) of the K{[}G{]}-module V. The space V is the direct sum of the family \((V^{\chi})_{\chi\in\widehat{\mathrm{G}}}\) , and the projector \(p_{\chi}\) of V with image \(V^{\chi}\) associated with this decomposition is given by

\[
p _ {\chi} = | \mathrm{G} | ^ {- 1} \sum_ {g \in \mathrm{G}} \chi (g ^ {- 1}) \pi (g)\tag{41}
\]

because of relation (19).

Remark. --- Let n be the cardinal of the group G, and let \(\mu_{n}(\mathrm{K})\) be the group of n-th roots of unity in K. For every \(g \in G\) , we have \(g^{n} = 1\) ; consequently, \(\widehat{G}\) can be identified with the group \(\operatorname{Hom}(\mathbf{G}, \mu_{n}(\mathbf{K}))\) . The group \(\mu_{n}(\mathbf{K})\) is cyclic of order n (V, §11, No.~2, p.~78, Theorem 1). The group \(\widehat{G}\) is therefore isomorphic to the group \(\mathrm{D}(\mathrm{G}) = \operatorname{Hom}(\mathrm{G}, \mathbf{Q}/\mathbf{Z})\) . By VII, §4, No.~9, p.~26,

Proposition 10, the group \(\widehat{\mathbf{G}}\) is isomorphic to the group \(\mathbf{G}\) , and the mapping that sends an element \(g\) of \(\mathbf{G}\) to the homomorphism \(\chi \mapsto \chi(g)\) from \(\widehat{\mathbf{G}}\) to \(\mathrm{K}^*\) is an isomorphism from \(\mathbf{G}\) to \(\widehat{\widehat{\mathbf{G}}}\) .
